\documentclass[12pt]{article}
\usepackage{amsmath,amssymb,amsthm}
\usepackage{geometry}
\usepackage{microtype}
\usepackage{hyperref}
\usepackage{natbib}

\geometry{margin=1.25in}

\newtheorem{theorem}{Theorem}[section]
\newtheorem{proposition}[theorem]{Proposition}
\newtheorem{corollary}[theorem]{Corollary}
\newtheorem{lemma}[theorem]{Lemma}
\newtheorem{definition}[theorem]{Definition}
\theoremstyle{remark}
\newtheorem{remark}[theorem]{Remark}

\title{Negation Before Logic:\\
Inferential Fields, Orientation Reversal,\\
and the Geometry of Admissible Inference}
\author{Flyxion\\
\small Independent Researcher}
\date{June 2026}

\begin{document}
\maketitle

\begin{abstract}
The cognitive cost of negation has been studied for decades, yet the 
dominant accounts remain unsatisfying. Operator-counting models predict 
that processing cost should increase monotonically with the number of 
negative expressions; they do not explain why two negations can be easier 
than one, why cost tracks monotonicity domains rather than operator 
occurrence, or why negated representations retain a processing signature 
long after the sentence has been parsed. Polarity-based accounts handle 
some scalar phenomena but conflate orientation toward zero with reversal 
of inferential direction, which are demonstrably distinct costs.

This paper proposes a geometric alternative. We introduce the notion of 
an \emph{inferential field} $\mathcal{F} = (X, R, \omega, A)$, where $X$ 
is a space of distinctions, $R$ is an inferential reachability relation, 
$\omega : X \to [-1,+1]$ is a continuous orientation field, and $A$ is an 
admissibility structure. Within this framework, negation is not 
fundamentally a truth-value operator. It is an orientation transformation 
$N : \mathcal{F} \to \mathcal{F}$ satisfying $N(\omega) = -\omega$. 
Monotonicity becomes a property of the orientation field rather than of 
individual lexical items. Upward-entailing environments correspond to 
regions where $\omega > 0$ and inferential reachability expands; 
downward-entailing environments correspond to regions where $\omega < 0$ 
and reachability contracts. Processing cost is not determined by 
operator count but by two separable quantities: the magnitude of 
orientation reversal and the curvature of admissibility domain boundaries.

This decomposition yields the \emph{Orientation-Reversal Principle}: the 
cognitive cost of a semantic operator is a function of the induced 
distortion of inferential orientation and the local curvature of 
admissibility domains, not of operator count. Four bodies of experimental 
evidence follow as consequences: the non-additivity of double negation, 
the domain-sensitivity of monotonicity processing cost, the sustained 
effect of negation in delayed verification, and the licensing conditions 
of negative polarity items. We further show that negative polarity item 
licensing is an admissibility probe: an expression such as \emph{ever} or 
\emph{any} succeeds not because a negative word is present but because 
the local region of the inferential field carries the required orientation 
and satisfies the admissibility condition jointly. Neural dissociation 
evidence from logical negation studies constrains the framework 
empirically, supporting a four-stage functional architecture in which 
syntactic field construction, logical field deformation, and scenario 
verification are partially separable operations.

The framework connects to a broader research program in admissibility 
geometry and distinguishability theory, in which the same mathematical 
structure --- distinctions, reachability, orientation, admissibility --- 
recurs across language, memory, repair, and scientific explanation. 
Negation, on this view, is not an addition to logic. It is a 
transformation of the distinction space within which logic becomes 
possible.
\end{abstract}

\newpage

\section{The Operator-Count Failure}

The simplest theory of negation processing is also the most natural one.
If a sentence contains a negative expression, it is harder to process than
the corresponding positive sentence. If it contains two negative
expressions, it is harder still. Processing cost increases monotonically
with the number of negative operators encountered during parsing. Call
this the \emph{operator-count hypothesis}.

The operator-count hypothesis has the virtue of simplicity and the vice
of being wrong in precisely the cases that matter most theoretically.
Three bodies of experimental evidence establish its failure, and the
pattern of failure is informative: it is not random noise around a
correct trend but a systematic deviation that points toward a different
underlying variable.

\subsection{The Double-Negation Asymmetry}

If processing cost tracked operator count, then a sentence containing two
negative expressions should be more costly than a sentence containing one.
The prediction is monotone: each additional negation adds a fixed
increment of difficulty. The experimental literature on double negation
does not support this prediction.

Sentences of the form \emph{it is not the case that none of the circles
are red} involve two scope-taking negative expressions and require the
reasoner to compute a positive conclusion by composing two reversals.
Under operator-count, this double-reversal configuration should be the
most expensive. Under an orientation account, the two reversals cancel:
$N(N(\omega)) = \omega$, returning the field to its original orientation.
The field does not count operators; it tracks net orientation. The
prediction from an orientation account is therefore that double negation
should not cost twice as much as single negation, and may in some
configurations cost less, because the inferential field has returned to a
familiar expansion regime.

This is what the experimental record shows. The cost of double negation
is not additive. In some participant populations and task configurations,
sentences with two negations are verified faster than sentences with one,
when syntactic complexity is controlled. Operator-count has no account of
this. Orientation reversal explains it directly.

\subsection{Domain Sensitivity of Monotonicity Cost}

A more precise failure of operator-count comes from experiments that
manipulate monotonicity independently of the number of negative
expressions. Consider the contrast between \emph{more than half} and
\emph{fewer than half}. Both are quantifiers. Neither is a negative word
in the lexical sense. Yet they induce different entailment profiles:
\emph{more than half} is upward-entailing in its restrictor, while
\emph{fewer than half} is downward-entailing. An operator-count model,
applied to these expressions, predicts no difference in processing cost,
since no negative operator has been added or removed.

The experimental evidence contradicts this prediction. Downward-entailing
quantifiers of this kind impose a measurable processing cost relative to
their upward-entailing counterparts, even when the sentences are matched
for length, frequency, syntactic structure, and the absence of any overt
negative morpheme. The source of the cost is the direction of inferential
expansion, not the presence of a negative word.

More decisive still is the finding from the 2024 monotonicity-domain
experiments. In those studies, the relevant experimental variable is not
the number of downward-entailing operators but the structure of the
monotonicity domain: whether a given expression occurs within a region
whose inferential orientation is uniformly downward, or near a boundary
where orientation shifts. The finding is that processing cost tracks
domain structure. Two configurations with equal operator counts but
different domain structures produce different costs. Two configurations
with unequal operator counts but equivalent domain structures produce
equivalent costs. This result cannot be accommodated by operator-count in
any form. It requires a variable that tracks the geometry of inferential
regions, not the inventory of lexical items.

\subsection{The Persistence of Negation in Delayed Verification}

A third failure of operator-count concerns the temporal profile of
negation effects. The operator-count hypothesis is naturally interpreted
as a parsing hypothesis: negative operators introduce additional
computational steps during incremental sentence processing, and those
steps consume time. On this interpretation, if sufficient processing time
is provided before a verification judgment is required, the cost should
diminish or disappear. The parser has had time to complete its work.

Delayed verification experiments test this prediction by introducing a
temporal gap between sentence presentation and the onset of the
verification image. The operator-count prediction is that negation effects
should attenuate with delay. The experimental finding is the opposite:
negation effects are \emph{sustained} across delays. A negated sentence
does not become easier to verify simply because more time has passed. The
cost persists into the verification phase even when parsing has long
concluded.

This result is incompatible with a purely incremental parsing account of
negation cost. It suggests that negation does not merely insert additional
steps into the parse; it changes the structure of the working-memory
representation that is passed forward to verification. In the framework
developed here, this is the \emph{persistence cost}: the residual
distortion of the inferential field that remains in the representation
after processing. Negation changes the orientation of the stored
representation, and orientation is not automatically restored by the
passage of time.

\subsection{What the Failures Share}

The three failures of operator-count share a common structure. In each
case, the quantity that predicts processing cost is not the number of
negative expressions but a property of the inferential environment those
expressions create: the net orientation of the field after all reversals
have been applied, the geometry of domain boundaries, and the persistence
of orientation change in the stored representation. These are properties
of a field, not properties of a count. The theory required to explain
them must be a field theory.

\section{Inferential Fields}

The failures documented in the preceding section converge on a single
diagnostic: the relevant variable is not a count of operators but a
property of the inferential environment those operators construct. To make
this precise, we need a mathematical object that can represent
inferential environments, support a notion of orientation, and admit a
well-defined transformation corresponding to negation. We introduce this
object here.

\subsection{Distinctions and Reachability}

Let $X$ be a set of \emph{distinctions}. Intuitively, a distinction is a
difference that the inferential system can represent and act upon: the
difference between red and not-red, between all and some, between the
circle being inside the square and outside it. We do not require $X$ to
carry propositional structure at the outset. Propositions, on the view
developed here, are derivative objects constructed from distinctions; the
geometry of the distinction space is prior.

\begin{definition}[Inferential reachability]
A binary relation $R \subseteq X \times X$ is an \emph{inferential
reachability relation} on $X$ if it satisfies the following conditions:
\begin{enumerate}
\item[(i)] \emph{Reflexivity}: $R(x,x)$ for all $x \in X$.
\item[(ii)] \emph{Transitivity}: if $R(x,y)$ and $R(y,z)$ then $R(x,z)$.
\item[(iii)] \emph{Non-triviality}: there exist $x, y \in X$ such that
$R(x,y)$ and $\neg R(y,x)$.
\end{enumerate}
\end{definition}

Condition (iii) ensures that $R$ is not the trivial equivalence relation.
Inferential motion has direction: some distinctions are reachable from
others without the converse holding. This directedness is what makes
orientation a meaningful concept in the space.

For a subset $A \subseteq X$, define the \emph{inferential expansion} of
$A$ under $R$ as

\[
R(A) = \{ y \in X : \exists\, x \in A,\; R(x,y) \}.
\]
This is the set of distinctions reachable from any member of $A$ in one
inferential step. The behavior of $R(A)$ as $A$ varies is what
monotonicity theory studies, and it is where orientation enters.

\subsection{The Orientation Field}

Classical monotonicity theory classifies expressions as upward- or
downward-entailing and treats this classification as a binary, lexically
determined property. The operator-count failures suggest that this
classification is insufficient: what matters is not the sign of an
individual expression but the orientation of the inferential environment
that expression inhabits. We represent this environment with a continuous
field.

\begin{definition}[Orientation field]
An \emph{orientation field} on $(X, R)$ is a continuous function
$\omega : X \to [-1, +1]$ satisfying:
\begin{enumerate}
\item[(i)] \emph{Upward regime}: if $\omega(x) > 0$ and $A \subseteq B
\subseteq X$, then $R(A) \subseteq R(B)$, so inferential expansion
preserves set inclusion.
\item[(ii)] \emph{Downward regime}: if $\omega(x) < 0$ and $A \subseteq
B \subseteq X$, then $R(B) \subseteq R(A)$, so inferential expansion
reverses set inclusion.
\item[(iii)] \emph{Boundary}: if $\omega(x) = 0$, then $x$ lies on the
boundary between upward and downward regimes, and the local expansion
behavior of $R$ is indeterminate.
\end{enumerate}
\end{definition}

The choice of $[-1,+1]$ rather than $\{-1,+1\}$ is deliberate. It allows
$\omega$ to vary continuously across the distinction space, so that
regions near orientation boundaries are represented as having intermediate
orientation values rather than being forced into a binary classification.
This continuity is what permits a well-defined notion of orientation
curvature in the next subsection.

\subsection{Admissibility}

Not every distinction in $X$ is available for inferential use at every
point in processing. Licensing conditions, syntactic constraints, and
contextual restrictions determine which regions of the distinction space
are open to a given inferential operation. We represent this with an
admissibility structure.

\begin{definition}[Admissibility structure]
An \emph{admissibility structure} on $(X, R, \omega)$ is a function
$A : X \to \{0, 1\}$ such that $A(x) = 1$ if the distinction $x$ is
currently admissible for inferential use and $A(x) = 0$ otherwise. An
\emph{admissible region} is a connected component of $A^{-1}(1)$.
\end{definition}

Admissibility is not truth. A distinction may be admissible without being
verified, and verified without being admissible in the relevant
inferential context. The separation between admissibility and truth is
essential to the account of negative polarity items developed in Section
4: an NPI is sensitive to admissibility conditions, not directly to truth
conditions.

\begin{definition}[Inferential field]
An \emph{inferential field} is a quadruple $\mathcal{F} = (X, R, \omega,
A)$ where $X$ is a set of distinctions, $R$ is an inferential
reachability relation on $X$, $\omega : X \to [-1,+1]$ is an orientation
field, and $A : X \to \{0,1\}$ is an admissibility structure.
\end{definition}

\subsection{Negation as Orientation Transformation}

With the inferential field in hand, we can give negation a precise
geometric meaning.

\begin{definition}[Negation operator]
The \emph{negation operator} $N : \mathcal{F} \to \mathcal{F}$ is the
map that acts on the orientation field by pointwise reversal:
\[
N(\omega)(x) = -\omega(x) \quad \text{for all } x \in X,
\]
while leaving $X$, $R$, and $A$ unchanged.
\end{definition}
Several consequences follow immediately.

\begin{proposition}[Double negation]
\label{prop:double-neg}
$N \circ N = \mathrm{id}_{\mathcal{F}}$.
\end{proposition}

\begin{proof}
$(N \circ N)(\omega)(x) = N(-\omega(x)) = -(-\omega(x)) = \omega(x)$ for
all $x \in X$.
\end{proof}

\begin{proposition}[Orientation of upward and downward environments]
\label{prop:orientation}
A linguistic environment is upward-entailing if and only if the
corresponding region of $\mathcal{F}$ satisfies $\omega > 0$. It is
downward-entailing if and only if the corresponding region satisfies
$\omega < 0$.
\end{proposition}

\begin{proof}
By conditions (i) and (ii) of Definition 2.2, the expansion behavior of
$R$ under set inclusion is determined by the sign of $\omega$. Upward
entailment is defined by the preservation of set inclusion under
inferential expansion; downward entailment by its reversal. These
correspond exactly to the upward and downward regimes of $\omega$.
\end{proof}

Proposition \ref{prop:double-neg} is already theoretically significant.
It establishes that the field-theoretic account of negation predicts the
cancelation of double negation as a structural theorem, not as an
empirical stipulation. The field does not count negative operators; it
tracks the net orientation that results from their composition. Two
reversals return the field to its original state. The non-additivity of
double-negation cost documented in the experimental literature is
therefore a consequence of the geometry, not a puzzle requiring separate
explanation.

\subsection{Orientation Curvature and Processing Cost}

The domain-sensitivity result from the 2024 monotonicity experiments
requires a further quantity: not merely the sign of $\omega$ at a point
but the rate at which $\omega$ changes across inferential regions. We
define this as orientation curvature.

\begin{definition}[Orientation curvature]
The \emph{orientation curvature} at $x \in X$ is
\[
\kappa_\omega(x) = \left| \nabla \omega(x) \right|,
\]
where $\nabla$ denotes the gradient with respect to the topology on $X$
induced by $R$.
\end{definition}
Intuitively, $\kappa_\omega(x)$ measures how rapidly the inferential
orientation changes in the neighborhood of $x$. In a uniformly
upward-entailing region, $\omega$ is constant and positive, so
$\kappa_\omega = 0$. In a uniformly downward-entailing region, $\omega$
is constant and negative, so $\kappa_\omega = 0$ again. High curvature
occurs at and near orientation boundaries, where $\omega$ transitions
between positive and negative values.

\begin{definition}[Processing cost]
The \emph{processing cost} at $x \in X$ is
\[
C(x) = \alpha \cdot \mathbf{1}[\omega(x) < 0] + \beta \cdot
\kappa_\omega(x),
\]
where $\alpha, \beta > 0$ are empirically determined weights,
$\mathbf{1}[\omega(x) < 0]$ is the indicator function for the downward
regime, and $\kappa_\omega(x)$ is orientation curvature.
\end{definition}
The first term captures the cost of operating in a reversed inferential
regime: downward-entailing environments impose a baseline cost because
the ordinary expansion direction of inference is no longer available. The
second term captures the additional cost of processing near domain
boundaries, where the local orientation is changing and neither a pure
upward nor a pure downward inference strategy is applicable.

This cost function separates three empirically distinct quantities that
previous accounts have conflated:

\begin{enumerate}
\item \emph{Polarity cost}: the cost associated with scalar negativity,
which in the present account is absorbed into the admissibility structure
$A$ rather than $\omega$ directly, since polarity is a property of the
scalar position of an expression rather than the orientation of its
inferential environment.

\item \emph{Monotonicity cost}: captured by the first term, the baseline
cost of operating in a downward-entailing region regardless of how that
region was constructed.

\item \emph{Boundary cost}: captured by the second term, the additional
cost of processing near orientation domain boundaries, which is the
quantity most directly measured by the 2024 monotonicity-domain
experiments.
\end{enumerate}

The persistence cost identified in delayed-verification studies does not
appear in $C(x)$ directly, because persistence is a property of the
representation passed forward to the verification stage rather than a
cost incurred during field navigation. We address persistence separately
in Section 5.

\subsection{Negation Distortion}

The final quantity introduced in this section connects the field-theoretic
account to the broader program in distinguishability geometry. In that
program, a transformation is characterized not only by its action on
objects but by the distortion it introduces in the reachability structure
of those objects. We define the analogous quantity for negation.

\begin{definition}[Negation distortion]
The \emph{negation distortion} at $x \in X$ is
\[
D_N(x) = d_R(x, N(x)),
\]
where $d_R$ is the reachability pseudometric on $X$ defined by
\[
d_R(x, y) = \inf \{ n \in \mathbb{N} : R^n(x) \ni y \text{ and }
R^n(y) \ni x \},
\]
with $d_R(x,y) = \infty$ if no such $n$ exists.
\end{definition}
Negation distortion measures how far a distinction $x$ is displaced in
the reachability structure of $X$ when the orientation of the field is
reversed. A distinction that remains easily reachable from its negated
counterpart has low distortion; a distinction that becomes unreachable or
reaches only after many inferential steps has high distortion.

The global distortion of the negation operator is then
\[
\|D_N\| = \int_X D_N(x)\, d\mu(x),
\]
where $\mu$ is a suitable measure on $X$ reflecting the prior probability
or salience of distinctions. This global quantity is what the overall
processing cost of negation in a given inferential context approximates.
The three cost components now have a unified representation: polarity
cost, monotonicity cost, and boundary cost all contribute to $\|D_N\|$
through their respective effects on $d_R$. Persistence cost is the
residual of $D_N$ that has not been resolved by the time the
representation is passed to the verification stage.

\section{The Orientation-Reversal Principle}

The formal machinery of Section 2 permits a precise statement of the
central theoretical claim of this paper. We state it as a principle
rather than a theorem because its full justification draws on both the
mathematical structure of inferential fields and the empirical record
reviewed in Section 1. The principle unifies the four experimental
phenomena --- double-negation non-additivity, domain-sensitive
monotonicity cost, sustained verification effects, and negative polarity
item licensing --- as consequences of a single underlying geometric
quantity.

\begin{theorem}[Orientation-Reversal Principle]
\label{thm:orp}
Let $\mathcal{F} = (X, R, \omega, A)$ be an inferential field and let
$T : \mathcal{F} \to \mathcal{F}$ be a semantic operator. The cognitive
processing cost of $T$ is determined not by the count of negative
expressions introduced by $T$ but by two separable quantities:
\begin{enumerate}
\item[(i)] the magnitude of net orientation reversal induced by $T$,
measured by $\int_X |\omega(x) - T(\omega)(x)|\, d\mu(x)$, and
\item[(ii)] the curvature of admissibility domain boundaries in the
region of $\mathcal{F}$ traversed during processing, measured by
$\int_X \kappa_\omega(x)\, d\mu(x)$ over the traversed region.
\end{enumerate}
\end{theorem}

\begin{proof}[Proof sketch]
We proceed by showing that operator count is neither necessary nor
sufficient to determine processing cost, and that quantities (i) and (ii)
are jointly sufficient within the field-theoretic framework.

\medskip
\noindent\textit{Operator count is not sufficient.} By the double-negation
result (Proposition \ref{prop:double-neg}), two orientation reversals
compose to the identity transformation. A sentence containing two
downward-entailing operators therefore induces a net orientation change of
zero, since $N \circ N = \mathrm{id}$. Under operator count, the
predicted cost is twice the single-negation cost. Under the
orientation-reversal measure (i), the predicted cost is zero for the
orientation component, plus boundary costs incurred during the traversal.
The experimental record supports the orientation prediction over the
operator-count prediction: double-negation cost is non-additive and
can fall below single-negation cost under controlled conditions.

\medskip
\noindent\textit{Operator count is not necessary.} By Proposition
\ref{prop:orientation}, downward-entailing environments correspond to
regions where $\omega < 0$, regardless of whether a negative morpheme is
present. Quantifiers such as \emph{fewer than half} induce downward
orientation without containing any negative word. The processing cost
associated with such quantifiers --- documented in the monotonicity
literature --- is therefore a cost of orientation, not a cost of operator
occurrence. Operator count, which registers zero for these expressions,
fails to predict the observed cost.

\medskip
\noindent\textit{Quantities (i) and (ii) are jointly sufficient.} The
cost function $C(x) = \alpha \cdot \mathbf{1}[\omega(x) < 0] + \beta
\cdot \kappa_\omega(x)$ captures the baseline cost of reversed orientation
and the additional cost of domain boundaries respectively. Integrating
over the traversed region yields a total cost that depends on the geometry
of the inferential path, not on the inventory of operators along it. This
accounts for: the non-additivity of double negation (net reversal is
zero); the domain-sensitivity of monotonicity cost (cost tracks boundary
curvature, not operator occurrence); and the differential cost of
downward-entailing quantifiers without negative morphemes (orientation
cost with zero operator count). The persistence effect requires an
additional claim about the representation passed to verification, which
is addressed in Section 5, but is consistent with the present framework
since $D_N(x) > 0$ whenever orientation has been reversed and the
distortion has not been resolved.
\end{proof}

\subsection{Four Phenomena as Consequences}

We now show explicitly how the four experimental phenomena documented in
Section 1 follow from the Orientation-Reversal Principle.

\subsubsection{Double-Negation Non-Additivity}

Under the operator-count hypothesis, the cost of a doubly negated
sentence $S^{--}$ satisfies $C(S^{--}) = 2 \cdot C(S^-)$, where $S^-$
is the singly negated sentence. Under the Orientation-Reversal Principle,
the relevant quantity is the net orientation after composition of
operators. Since $N \circ N = \mathrm{id}$, the net reversal induced by
$S^{--}$ is zero, and the orientation component of cost vanishes. The
residual cost is purely the boundary cost incurred during the processing
of the intermediate state, which lies at the orientation boundary
$\omega = 0$. This boundary cost is real --- the system does pass through
a region of high curvature --- but it is bounded above by the single
boundary traversal, not multiplied by two.

The empirical prediction is therefore:
\[
C(S^{--}) < 2 \cdot C(S^-),
\]
and in populations or task configurations where boundary traversal cost
is low relative to orientation cost, the stronger prediction
\[
C(S^{--}) < C(S^-)
\]
may hold. Both predictions are supported by the experimental record and
both are consequences of the geometry.
\subsubsection{Domain-Sensitive Monotonicity Cost}

The 2024 monotonicity-domain experiments show that processing cost tracks
domain structure rather than operator occurrence. Under the
Orientation-Reversal Principle, this is expected. The cost function
$C(x)$ depends on $\omega(x)$ and $\kappa_\omega(x)$, both of which are
properties of the inferential region in which processing occurs, not of
the count of operators that constructed that region. Two configurations
with equal operator counts but different domain structures will have
different values of $\kappa_\omega$ and therefore different costs. Two
configurations with unequal operator counts but equivalent domain
structures will have equivalent costs.

More precisely, consider two sentences $S_1$ and $S_2$ such that:
\begin{itemize}
\item $S_1$ contains one downward-entailing operator and processing
occurs uniformly within a single downward-entailing domain.
\item $S_2$ contains two downward-entailing operators but processing
traverses a boundary between an upward and a downward domain.
\end{itemize}
Operator count predicts $C(S_2) > C(S_1)$. The Orientation-Reversal
Principle predicts $C(S_2) > C(S_1)$ only if the boundary cost of $S_2$
exceeds the difference in baseline orientation cost, which depends on the
specific geometry of the domains. The empirical result --- that $S_2$ can
be more costly than predicted by its operator count and $S_1$ less costly
--- follows naturally from the curvature term $\beta \cdot
\kappa_\omega(x)$.

\subsubsection{Sustained Verification Effects}

The delayed-verification result presents a challenge to any account that
treats negation cost as a purely incremental parsing phenomenon. Under
the Orientation-Reversal Principle, the sustained effect has a natural
explanation via negation distortion. When a negated sentence is
processed, the orientation of the resulting representation is reversed
relative to the corresponding positive sentence. This reversed
representation is stored in working memory and passed to the verification
stage. The distortion $D_N(x) = d_R(x, N(x))$ quantifies how far the
representation has been displaced in reachability terms. This distortion
does not automatically resolve with the passage of time; it persists in
the representation until the verification process has explicitly
re-oriented the field with respect to the scenario being compared.

The prediction is that the sustained cost should correlate with the
magnitude of $D_N$ for the relevant distinctions, and that it should
diminish only when the verification process has completed the reachability
comparison --- not merely when parsing time has elapsed.

\subsubsection{Negative Polarity Item Licensing}

The analysis of negative polarity item licensing requires the admissibility
structure $A$ of the inferential field and is developed in full in Section
4. Here we note only the consequence of the Orientation-Reversal
Principle that is relevant to licensing: an NPI is not sensitive to the
presence of a negative operator but to the orientation and admissibility
of the local inferential region. Two environments that are equally
negative in operator count but differ in orientation or admissibility will
differ in NPI licensing. This prediction separates the present account
from both standard downward-entailment accounts, which license NPIs by
operator type, and from polarity accounts, which license them by scalar
position.

\subsection{The Principle as Unification}

The Orientation-Reversal Principle is not merely a restatement of the
four phenomena in geometric language. It provides a single quantity ---
the distortion of inferential orientation and its curvature across
admissibility domains --- from which all four phenomena are derivable.
This is a stronger claim than a descriptive unification. It implies that
the four phenomena are not independently explained by four separate
mechanisms but are four measurements of the same underlying geometric
variable. Experiments designed to manipulate one of the four phenomena
while holding the geometric variable constant should attenuate or
eliminate the effect. Experiments designed to manipulate the geometric
variable directly --- by constructing sentences that traverse domain
boundaries without using negative morphemes, for instance --- should
produce costs not predicted by any account that relies on negative
operator occurrence.

These are testable predictions. We return to them in Section 6.

\section{Negative Polarity Items as Admissibility Probes}

Negative polarity items present the licensing problem in its sharpest
form. An expression such as \emph{ever}, \emph{any}, \emph{yet}, or
\emph{lift a finger} is grammatical in some environments and
ungrammatical in others, and the distribution of grammatical environments
does not reduce to the mere presence of a negative word. The sentences in
(\ref{ex:npi-basic}) illustrate the basic pattern.

\begin{enumerate}
\item \label{ex:npi-basic}
\begin{enumerate}
\item[(a)] Nobody has ever been to that island. \hfill [grammatical]
\item[(b)] Somebody has ever been to that island. \hfill
[ungrammatical]
\item[(c)] Few students have ever read that book. \hfill [grammatical]
\item[(d)] Many students have ever read that book. \hfill
[ungrammatical]
\item[(e)] I doubt that she has ever been there. \hfill [grammatical]
\item[(f)] I believe that she has ever been there. \hfill
[ungrammatical]
\end{enumerate}
\end{enumerate}

The standard generalization, due to Ladusaw, is that NPIs are licensed
in downward-entailing environments. \emph{Nobody}, \emph{few}, and the
complement of \emph{doubt} are downward-entailing; \emph{somebody},
\emph{many}, and the complement of \emph{believe} are not. This
generalization is largely correct as a descriptive matter, but it leaves
two theoretical questions unanswered. First, why should downward
entailment be the licensing condition? What is the connection between
inference reversal and the grammaticality of \emph{ever}? Second, how
does the licensing condition interact with syntactic scope and domain
structure, given the 2024 finding that processing cost tracks domains
rather than operators?

The inferential field framework answers both questions by recasting NPI
licensing as an admissibility probe.

\subsection{Admissibility and Orientation as Independent Conditions}

The key theoretical move is to separate two conditions that previous
accounts have run together. In the standard downward-entailment account,
a single condition --- the DE-ness of the licensing environment ---
determines grammaticality. In the field-theoretic account, two conditions
must be jointly satisfied.

\begin{definition}[NPI probe]
An \emph{NPI probe} is a function $P : X \to \{0, 1\}$ defined by
\[
P(x) = \begin{cases} 1 & \text{if } A(x) = 1 \text{ and } \omega(x) <
0, \\ 0 & \text{otherwise.} \end{cases}
\]
An NPI expression is licensed at position $x$ in an inferential field
$\mathcal{F} = (X, R, \omega, A)$ if and only if $P(x) = 1$.
\end{definition}
The two conditions are:
\begin{enumerate}
\item[(i)] \emph{Admissibility}: $A(x) = 1$. The distinction at $x$ must
be open to inferential use in the current context. This is a structural
condition on the inferential field, sensitive to syntactic scope,
embedding depth, and contextual licensing factors.
\item[(ii)] \emph{Orientation}: $\omega(x) < 0$. The local inferential
environment must be in a downward regime. This is the orientation
condition, which corresponds to the DE-ness condition in standard
accounts but is now understood as a property of the field rather than
of a lexical item.
\end{enumerate}

These conditions are genuinely independent. An environment can be
downward-oriented without being admissible for NPI licensing --- certain
embedded contexts are downward-entailing but block NPI licensing for
independent structural reasons. Conversely, an environment can be
admissible in general while failing to carry the required downward
orientation. NPI licensing fails in both cases, but for different
reasons, and the distinction between the two failure modes has empirical
consequences.

\subsection{Why Downward Orientation Licenses NPIs}

The independence of the two conditions raises the explanatory question
more sharply: given that both conditions must be satisfied, what is the
theoretical reason that downward orientation is one of them? The
field-theoretic account provides an answer that the standard DE account
does not.

In an upward-entailing environment, inferential expansion proceeds
outward: from narrower to broader sets, from more specific to more
general distinctions. An existential expression such as \emph{ever} or
\emph{any} ranges over a class of distinctions and asserts the
non-emptiness of that class. In an upward-entailing environment, this
assertion is redundant with ordinary scalar inference: if the class is
non-empty at a specific level, upward inference carries the assertion to
more general levels automatically. The NPI contributes no additional
inferential work.

In a downward-entailing environment, by contrast, the orientation of
the field is reversed. Inference no longer expands outward; it contracts
inward. In this regime, the existence of an instance at a general level
does not guarantee existence at a specific level. The NPI \emph{ever}
now does genuine inferential work: it asserts reachability across a
region of the field where reachability is not guaranteed by the default
expansion of upward inference. The NPI is licensed precisely because the
field orientation creates a gap that the NPI's assertion is needed to
fill.

This gives a functional explanation of the licensing condition. NPIs are
not licensed by downward entailment as an arbitrary syntactic trigger.
They are licensed because downward orientation creates the inferential
conditions under which their semantic contribution is non-trivial. An
NPI in an upward-entailing environment contributes nothing that ordinary
inference does not already deliver; an NPI in a downward-entailing
environment contributes an assertion that the field orientation makes
genuinely informative.

\begin{proposition}[Licensing and inferential contribution]
\label{prop:npi-contrib}
Let $\mathcal{F} = (X, R, \omega, A)$ be an inferential field and let
$e$ be an NPI expression at position $x$. The inferential contribution
of $e$ is non-trivial if and only if $\omega(x) < 0$.
\end{proposition}

\begin{proof}
Suppose $\omega(x) > 0$. Then the local region is upward-entailing, and
by condition (i) of Definition 2.2, $R(A) \subseteq R(B)$ whenever $A
\subseteq B$. For an existential NPI ranging over a class $\mathcal{C}$
of distinctions, the assertion that $\mathcal{C}$ is non-empty at some
level is carried upward automatically: $R(\mathcal{C}) \supseteq
\mathcal{C}$, and the reachability of members of $\mathcal{C}$ is
entailed by the upward expansion of $R$. The NPI asserts what the field
already delivers. Its contribution is trivial.

Suppose $\omega(x) < 0$. Then the local region is downward-entailing,
and by condition (ii) of Definition 2.2, $R(B) \subseteq R(A)$ whenever
$A \subseteq B$. The default direction of inferential expansion is
reversed. The assertion that $\mathcal{C}$ is non-empty at some level
is no longer carried automatically: existence at a general level does
not guarantee reachability at a specific level under contraction.
The NPI must assert reachability explicitly. Its contribution is
non-trivial.
\end{proof}

\subsection{Domain Structure and Licensing Domains}

The 2024 monotonicity-domain finding that processing cost tracks domain
structure rather than operator count has a direct parallel in NPI
licensing. If licensing were purely a property of individual operators,
then the licensing condition would be checked locally at each operator
without reference to the broader domain structure. The field-theoretic
account predicts otherwise: licensing is a property of admissible
orientation regions, and those regions have spatial extent.

\begin{definition}[Licensing domain]
A \emph{licensing domain} for NPIs in $\mathcal{F}$ is a maximal
connected component $D \subseteq X$ such that $A(x) = 1$ and $\omega(x)
< 0$ for all $x \in D$.
\end{definition}

Several predictions follow from this definition.

First, an NPI is licensed throughout a licensing domain, not merely at
the position of the licensing operator. The domain has spatial extent,
and the orientation condition is satisfied everywhere within it. This
explains the well-known fact that NPIs can be licensed at a distance from
their licensor, within the bounds of the appropriate syntactic scope.

Second, the boundary of a licensing domain is a region of high
orientation curvature: $\kappa_\omega$ is large near the boundary. The
processing cost function $C(x) = \alpha \cdot \mathbf{1}[\omega(x) < 0]
+ \beta \cdot \kappa_\omega(x)$ therefore predicts that sentences where
an NPI appears near the edge of a licensing domain --- where the field
transitions from downward to upward orientation --- should be more
costly to process than sentences where the NPI appears well within the
domain. This is a testable prediction that distinguishes the field
account from standard licensing accounts.

Third, an NPI that appears outside a licensing domain but within an
admissible region ($A(x) = 1$, $\omega(x) > 0$) should produce an
ungrammaticality judgment associated with orientation mismatch rather
than admissibility failure. The two failure modes predict different
patterns of degradation and potentially different processing signatures.
An NPI that appears outside an admissible region altogether ($A(x) = 0$)
should produce a different kind of ungrammaticality. Whether these
predicted differences in degradation type correspond to empirically
distinguishable judgment profiles is an open question.

\subsection{Superpolarity and Weak NPIs}

The standard DE account treats all NPIs as requiring the same licensing
condition. The empirical literature, however, distinguishes weak NPIs
(licensed in a wide range of DE environments, including questions and
conditionals) from strong NPIs (licensed only in a narrow range of
environments, typically clausemate negation). This distinction is
difficult to capture in a purely operator-based account.

The field-theoretic account handles the distinction through the joint
conditions of admissibility and orientation. A weak NPI requires only
that $\omega(x) < 0$; it is insensitive to the admissibility structure
except in the minimal sense that the position must be syntactically
accessible. A strong NPI imposes an additional admissibility condition:
it requires not merely that the local orientation be downward but that
the admissibility structure satisfy a further constraint, which in
practice corresponds to the requirement that the licensor be a clausemate
negative operator rather than a more distant downward-entailing
expression.

\begin{definition}[Weak and strong NPI probes]
A \emph{weak NPI probe} at $x$ succeeds if $\omega(x) < 0$, regardless
of the admissibility structure beyond syntactic accessibility. A
\emph{strong NPI probe} at $x$ succeeds if $\omega(x) < -\theta$ for
some threshold $\theta \in (0,1)$ and $A(x) = 1$ under a restricted
admissibility structure $A_s$ that requires clausemate licensing.
\end{definition}

The threshold $\theta$ in the strong NPI definition captures the
intuition that strong NPIs require not merely a downward-oriented
environment but a strongly downward-oriented one: the orientation must
exceed a threshold magnitude, not merely be negative. This is a natural
consequence of allowing $\omega$ to range continuously over $[-1,+1]$
rather than taking binary values. Environments near the boundary
$\omega(x) = 0$ are weakly downward-oriented and license weak but not
strong NPIs; environments deep in the downward regime license both.

\subsection{The Probe Metaphor and Its Limits}

We have described NPIs as probes of the inferential field, but the
metaphor should not be pressed too far. An NPI is not an independent
detector inserted into the sentence to measure field properties. It is a
lexical item whose semantic contribution is constituted by its
sensitivity to field orientation: it asserts reachability across a
downward-oriented region, and its grammaticality is the reflection of
whether such an assertion is coherent in the local environment. The probe
metaphor captures the epistemic situation of the linguist or the grammar:
by observing where NPIs are licensed, one can infer the orientation
structure of the inferential field. The NPI is a probe in the sense that
a thermometer is a probe: it does not measure temperature independently
of being a thermometer; its readings constitute evidence about the
environment precisely because of what it is.

The field-theoretic account therefore does not add NPIs as an external
diagnostic layer on top of the semantic theory. It integrates them as
lexical items whose licensing conditions are determined by the same
geometric structure that determines processing cost, monotonicity
behavior, and negation distortion. The grammar, the processing system,
and the geometric framework are not three separate theories coordinated
by stipulation. They are three descriptions of the same inferential field.

\section{Persistence, Memory, and Residual Distortion}

The delayed-verification experiments occupy a distinctive position in the
empirical record. The other phenomena discussed in this paper ---
double-negation non-additivity, domain-sensitive monotonicity cost, NPI
licensing --- concern the processing of sentences at or near the time of
their production. The delayed-verification result concerns what happens
after processing has nominally concluded. A negated sentence retains a
processing cost advantage for the corresponding positive sentence even
when a temporal gap is introduced between the sentence and the
verification image. The cost is not a parsing artifact that dissipates
with time; it is a property of the stored representation.

This result is significant for the field-theoretic account because it
establishes that the orientation transformation induced by negation is not
merely a transient computational state. It is a change in the structure
of the representation that persists in working memory and affects
downstream processing. We formalize this persistence using the concept of
residual distortion introduced in Section 2 and develop its consequences
for the architecture of the verification process.

\subsection{The Representation Passed to Verification}

In the four-stage model introduced in Section 2 and elaborated in Section
6, the verification stage receives as input the output of field
deformation: a representation whose orientation has been modified by the
logical operators in the sentence. For a positive sentence $S^+$, the
representation passed to verification is $C_\sigma(S^+)$, the field
constructed by syntactic processing. For a negative sentence $S^-$, the
representation passed to verification is $D_N(C_\sigma(S^-))$, the field
after orientation reversal.

The crucial observation is that this representation is not the same
object as $C_\sigma(S^+)$ with a flag attached. The orientation reversal
is not a notational marker that the verification process can simply read
off and invert. It is a structural change in the reachability geometry of
the representation: the distinctions that were mutually reachable in the
original field are no longer mutually reachable in the same pattern after
negation, and conversely. The verification process must navigate this
restructured geometry in order to compare the representation against the
visual scenario.

\begin{definition}[Residual distortion]
Let $\mathcal{F} = (X, R, \omega, A)$ be an inferential field and let
$\mathcal{F}^- = D_N(\mathcal{F}) = (X, R, -\omega, A)$ be the field
after orientation reversal. The \emph{residual distortion} of
$\mathcal{F}^-$ with respect to $\mathcal{F}$ at position $x$ is
\[
\Delta(x) = d_R(x, N(x)) = D_N(x),
\]
where $d_R$ is the reachability pseudometric of Definition 2.7. The
\emph{global residual distortion} is
\[
\|\Delta\| = \int_X \Delta(x)\, d\mu(x).
\]
\end{definition}
The residual distortion $\|\Delta\|$ is the quantity that the
delayed-verification result measures. It is not the cost of parsing the
sentence; that cost has already been incurred. It is the cost imposed on
the verification process by the fact that the representation it receives
has a different reachability geometry from the representation that would
have been passed by the corresponding positive sentence.

\subsection{Why Time Does Not Resolve Distortion}

The operator-count hypothesis, interpreted as a parsing hypothesis,
predicts that negation cost should diminish with delay because the
parsing operations that incur the cost have time to complete. The
field-theoretic account predicts the opposite: residual distortion
persists because it is a structural property of the stored representation,
not a measure of incomplete processing.

\begin{proposition}[Persistence of residual distortion]
\label{prop:persistence}
Let $\tau$ be a temporal delay between sentence offset and verification
image onset. Under the field-theoretic account, the processing cost of
verification for a negated sentence satisfies
\[
C_V(S^-, \tau) \geq C_V(S^+, \tau) + \|\Delta\|_\tau
\]
for all $\tau \geq 0$, where $\|\Delta\|_\tau$ is the residual
distortion at delay $\tau$, and $\|\Delta\|_\tau$ does not converge to
zero as $\tau \to \infty$ in the absence of active re-orientation.
\end{proposition}
\begin{proof}
The verification process requires computing $\chi(P, \mathcal{F}^-)$,
where $P$ is the visual scenario and $\mathcal{F}^-$ is the negated
representation. The cost of $\chi$ depends on the reachability distance
between the scenario representation and the sentence representation. For
a positive sentence, this distance is $d_R(P, \mathcal{F}^+)$. For a
negative sentence, the representation has been distorted by $D_N$, so
the effective distance is $d_R(P, \mathcal{F}^-)$. The difference
between these distances is bounded below by $\|\Delta\|$ by the triangle
inequality applied to $d_R$.

The temporal delay $\tau$ does not reduce $\|\Delta\|_\tau$ because
$\|\Delta\|$ is a property of the representation stored in working
memory, not a property of an ongoing computational process. The
representation does not spontaneously re-orient itself with the passage
of time; orientation is a structural feature of the stored field, not a
transient state of a processing operation. Therefore $\|\Delta\|_\tau$
remains bounded away from zero for all $\tau$, and the cost inequality
is maintained across all delays.

Active re-orientation --- the explicit computation of $N(\mathcal{F}^-)
= \mathcal{F}^+$ as a preliminary step before verification --- would
resolve the distortion, but this operation itself incurs a cost
proportional to $\|\Delta\|$. The total verification cost for a negative
sentence is therefore at least $C_V(S^+, \tau) + \|\Delta\|$ regardless
of delay, with the cost allocated either to working-memory maintenance or
to explicit re-orientation depending on strategy.
\end{proof}

\subsection{Two Verification Strategies}

The proposition above implies that reasoners facing a delayed
verification task with negated sentences have two available strategies,
both of which incur a cost related to $\|\Delta\|$.

The first strategy is \emph{deferred re-orientation}: maintain the
negated representation $\mathcal{F}^-$ in working memory throughout the
delay interval, and re-orient the field explicitly when the verification
image appears. The working-memory cost is the maintenance cost of storing
a distorted representation; the re-orientation cost is incurred at
verification onset. This strategy predicts that reaction time should
spike at the onset of the verification image relative to positive
sentences, with the spike proportional to $\|\Delta\|$.

The second strategy is \emph{anticipatory re-orientation}: re-orient the
field during the delay interval, converting $\mathcal{F}^-$ to
$\mathcal{F}^+$ before the image appears. This strategy amortizes the
re-orientation cost across the delay interval rather than concentrating
it at verification onset. It predicts a more uniform cost increase across
the delay interval rather than a spike at onset.

Both strategies predict a sustained cost relative to positive sentences.
They differ in the temporal distribution of that cost. The delayed-
verification experiments can in principle distinguish between them by
examining the time course of the reaction time disadvantage: a spike at
onset supports deferred re-orientation; a uniform elevation supports
anticipatory re-orientation. The available evidence is more consistent
with the deferred strategy, suggesting that working memory maintains the
distorted representation rather than re-orienting proactively, but the
evidence is not yet decisive.

\subsection{Persistence and Working-Memory Architecture}

The persistence result has implications for the architecture of working
memory that go beyond the specific case of negation. If residual
distortion is a genuine structural feature of stored representations,
then working memory must be capable of maintaining representations with
non-default orientation. This is a non-trivial architectural requirement:
it implies that working memory does not normalize representations to a
canonical positive orientation during maintenance.

This requirement is consistent with the broader evidence on working
memory and logical form. Working memory representations retain structural
properties of the sentences that generated them, including scope
relations, embedding depth, and logical structure. Orientation, in the
present framework, is one such structural property. Its maintenance in
working memory is therefore not a special feature of negation but an
instance of the general capacity to maintain structured logical
representations across delays.

What is special about negation is the direction of the distortion.
Upward-entailing sentences generate representations whose orientation
matches the default expansion direction of the inferential field.
Downward-entailing sentences generate representations whose orientation
is reversed. The reversed orientation imposes a greater maintenance cost
because it conflicts with the default reachability geometry that the
verification process is designed to navigate. The verification process
is tuned to a field in which $\omega > 0$; a representation with $\omega
< 0$ requires either re-orientation or navigation in the non-default
direction, both of which incur additional cost.

\begin{corollary}[Asymmetry of maintenance cost]
\label{cor:maintenance}
Working-memory maintenance of a negated representation $\mathcal{F}^-$
incurs a greater cost than maintenance of the corresponding positive
representation $\mathcal{F}^+$, proportional to $\|\Delta\|$.
\end{corollary}

\begin{proof}
Maintenance cost is the cost of preserving the structural properties of
a representation against decay and interference over the delay interval.
For $\mathcal{F}^+$, the orientation $\omega > 0$ is consistent with
the default expansion direction of $R$, so maintenance requires only
that the reachability structure be preserved. For $\mathcal{F}^-$, the
orientation $\omega < 0$ conflicts with the default expansion direction,
so maintenance requires additionally that the reversed orientation be
preserved against spontaneous re-normalization toward the default. This
additional maintenance requirement is proportional to $\|\Delta\|$,
since larger distortions require greater active maintenance to prevent
reversion to default orientation.
\end{proof}

\subsection{Persistence and the Admissibility Structure}

A further consequence of the persistence result concerns the
admissibility structure $A$ of the stored representation. We have
defined admissibility as a function on the distinction space that
determines which distinctions are open to inferential use. When a
negated representation is stored in working memory, the admissibility
structure is maintained alongside the orientation reversal.

This has an important consequence for NPI licensing in embedded
environments. An NPI that is licensed at the time of sentence processing
--- because $A(x) = 1$ and $\omega(x) < 0$ in the processed field ---
remains licensed in the stored representation, since both conditions are
preserved under maintenance. But the verification process may require
comparing the stored representation against a scenario in which the
admissibility conditions differ. In such cases, the comparison must
resolve a mismatch not only in orientation but in admissibility
structure. This mismatch adds an additional component to the verification
cost beyond the residual distortion $\|\Delta\|$.

This prediction extends the delayed-verification paradigm in a new
direction. Sentences containing NPIs in downward-entailing environments
should show a specific pattern: the NPI licensing condition is satisfied
at processing time, but the verification cost includes both the
orientation distortion cost and an admissibility mismatch cost if the
verification scenario does not independently support the admissibility
conditions. Designing verification experiments that manipulate the
admissibility match between sentence and scenario would test this
prediction directly.

\subsection{Summary: Three Components of Negation Cost}

The analysis of this section completes the decomposition of negation cost
into three separable components, each with a distinct temporal profile
and a distinct locus in the processing architecture.

\begin{enumerate}
\item \emph{Orientation cost}: incurred during field construction and
deformation, proportional to $\mathbf{1}[\omega(x) < 0]$. This is the
baseline cost of operating in a downward-entailing regime and is
concentrated in the processing stage.

\item \emph{Boundary cost}: incurred during traversal of orientation
domain boundaries, proportional to $\kappa_\omega(x)$. This cost is
concentrated near domain edges and is the quantity most directly
measured by the 2024 monotonicity-domain experiments.

\item \emph{Residual distortion cost}: incurred during working-memory
maintenance and verification, proportional to $\|\Delta\|$. This cost
persists across temporal delays and is the quantity measured by
delayed-verification experiments. It does not diminish with parsing
time but only with active re-orientation.
\end{enumerate}

These three components are in principle independently manipulable. A
sentence that traverses many domain boundaries may have high boundary
cost with moderate residual distortion. A sentence with a single deep
negation may have low boundary cost with high residual distortion. A
sentence with two canceling negations may have low orientation cost, high
boundary cost at the intermediate state, and near-zero residual
distortion. The three-component decomposition generates a matrix of
predictions across sentence types and task conditions that jointly
constrain the framework and distinguish it from both operator-count and
polarity accounts.

\section{Neural Dissociation and the Four-Stage Model}

The inferential-field account developed in the preceding sections is a
theoretical framework, not a neural model. It makes claims about the
computational structure of negation processing --- the geometry of
orientation reversal, the separability of cost components, the
persistence of residual distortion --- without committing to specific
neural implementations of those computations. Nevertheless, the framework
is not neurally unconstrained. It requires, as a necessary condition,
that the operations it distinguishes be at least partially dissociable at
the level of neural implementation. If syntactic field construction,
logical field deformation, and scenario verification were realized by a
single undifferentiated neural process, then the cost decomposition
developed here would be computationally real but neurally uninstantiated,
and the framework would lose its claim to describe the actual processing
architecture rather than an idealized one.

The neural evidence on negation processing provides the required
constraint. We review the relevant findings, formalize the dissociation
argument, and then use it to ground the four-stage functional model that
has been implicit throughout the paper.

\subsection{The Relevant Neural Evidence}

Three clusters of neural findings bear on the dissociation argument.

The first concerns syntactic complexity. A substantial literature
associates complex syntactic processing with left inferior frontal regions,
particularly Broca's area (BA 44 and BA 45) and adjacent premotor
cortex, together with posterior superior temporal regions. These regions
show increased activity for sentences with greater syntactic complexity
--- center-embedded relative clauses, non-canonical argument structures,
long-distance dependencies --- relative to structurally simpler sentences
matched for semantic content. The relevant point for present purposes is
not the precise anatomical attribution but the existence of a reliable
syntactic complexity contrast that localizes to identifiable regions.

The second concerns logical negation specifically. Grodzinsky and
colleagues report a left anterior insula cluster, overlapping with the
cytoarchitectonically defined region Id7, that exhibits negation-sensitive
activity in controlled sentence-picture verification tasks. The
experimental design isolates logical negation by contrasting negative
sentences against positive sentences matched for syntactic complexity,
word frequency, and sentence length. The negation-sensitive contrast
localizes to the anterior insula cluster rather than to the classical
syntactic regions implicated by the first set of findings. Critically,
participant-level reaction times in the negation condition correlate with
BOLD activity in this cluster, establishing a link between the neural
signal and the behavioral cost.

The third concerns numerical and scenario comparison. Tasks requiring
magnitude comparison or verification against a structured visual scenario
implicate a distinct set of regions, including intraparietal sulcus and
inferior parietal cortex, that are separable from both the syntactic and
the negation clusters. The verification stage of the sentence-picture
task is not neurally identical to the syntactic or logical processing
stages.

These three clusters --- syntactic regions, anterior insula negation
cluster, parietal comparison regions --- are anatomically distinct and
are differentially sensitive to distinct experimental contrasts. This
triple dissociation is the neural evidence that the framework requires.

\subsection{Formalizing the Dissociation}

We formalize the dissociation argument using the four-map decomposition
introduced in Section 2. A sentence-picture verification trial is
represented as the composite operation
\[
V(P, S) = \chi\bigl(P,\; D_\ell(C_{\sigma}(S))\bigr),
\]
where $S$ is the sentence, $P$ is the visual scenario, $C_{\sigma}$ is
syntactic field construction, $D_\ell$ is logical field deformation, and
$\chi$ is scenario comparison. For a positive sentence,
\[
V(P, S^+) = \chi(P, C_{\sigma}(S^+)),
\]
and for a negative sentence,
\[
V(P, S^-) = \chi(P, D_N(C_{\sigma}(S^-))).
\]
The contrast between negative and positive conditions isolates $D_N$
under the assumption that $C_{\sigma}$ and $\chi$ are held sufficiently
constant across conditions by the experimental design.
\begin{proposition}[Neural dissociation of field deformation]
\label{prop:dissociation}
Suppose the verification task decomposes as $V = \chi \circ D_\ell \circ
C_{\sigma}$. Let $R_C$, $R_D$, and $R_\chi$ denote neural regions
associated with syntactic construction, logical deformation, and scenario
comparison respectively. If
\begin{enumerate}
\item[(i)] the negation contrast modulates $B_{R_D}$ but not $B_{R_C}$
or $B_{R_\chi}$,
\item[(ii)] syntactic complexity manipulations modulate $B_{R_C}$ but
not $B_{R_D}$ or $B_{R_\chi}$, and
\item[(iii)] comparison difficulty manipulations modulate $B_{R_\chi}$
but not $B_{R_C}$ or $B_{R_D}$,
\end{enumerate}
then $C_{\sigma}$, $D_\ell$, and $\chi$ are functionally dissociable
operations.
\end{proposition}

\begin{proof}
Let $B_R(O)$ denote the BOLD response in region $R$ associated with
operation $O$. If the three operations were not dissociable --- if they
constituted a single undifferentiated process --- then any manipulation
that increases task difficulty should modulate the same neural network
regardless of whether the manipulation targets syntactic structure,
logical form, or comparison. The observed pattern is the opposite: each
manipulation produces a selective effect in its associated region while
leaving the other regions unaffected. This selectivity is
contrast-relative: the negation contrast is not simply a harder condition
in general, but a condition that specifically increases demand on $R_D$
while syntactic and comparison demands are controlled. The triple
dissociation therefore establishes that the three operations are
functionally separable, in the sense that each can be selectively
modulated without equivalent modulation of the others. Under the
decomposition $V = \chi \circ D_\ell \circ C_{\sigma}$, this selectivity
is the neural signature of functional dissociation.
\end{proof}

\subsection{Against the Generic Load Objection}

The anterior insula is implicated in a wide range of non-linguistic
functions, including interoception, affective processing, salience
detection, and cognitive control. The observation that a negation
contrast localizes partly to this region is therefore consistent with the
hypothesis that negation is simply harder in a generic sense, and that
the anterior insula is responding to generic task difficulty rather than
to the logical structure of the sentence. We address this objection
directly.

\begin{proposition}[Specificity over generic load]
\label{prop:specificity}
If anterior insula activity in negation tasks were caused by generic task
load rather than by logical field deformation, then the same region
should be modulated by any manipulation that increases task difficulty,
regardless of whether the manipulation targets syntactic structure,
logical form, or comparison. The observed pattern of anatomical
selectivity is inconsistent with this prediction.
\end{proposition}

\begin{proof}
Let $L$ be a scalar measure of generic task load. On the generic-load
hypothesis, regional activity is a monotone function of $L$:
\[
B_R = f_R(L), \qquad f_R' \geq 0.
\]
Under this hypothesis, any manipulation that increases $L$ should
increase $B_{R_D}$, $B_{R_C}$, and $B_{R_\chi}$ proportionally, since
all reflect the same underlying load variable. The predicted pattern is
uniform modulation across regions for any difficulty-increasing
manipulation.
The observed pattern is not uniform. Syntactic complexity increases
$B_{R_C}$ without proportionally increasing $B_{R_D}$ or $B_{R_\chi}$.
Negation increases $B_{R_D}$ without proportionally increasing $B_{R_C}$
or $B_{R_\chi}$. Comparison difficulty increases $B_{R_\chi}$ without
proportionally increasing the other regions. The generic-load function
$f_R(L)$ cannot produce this pattern, since it predicts that any common
increase in $L$ should modulate all load-sensitive regions.

The residual explanation available to the generic-load account is that
the three regions have different sensitivity functions $f_{R_C}$,
$f_{R_D}$, $f_{R_\chi}$ with different thresholds or slopes, and that
the three manipulations happen to produce load increases that fall
selectively above each region's threshold. This explanation is available
in principle but requires three independent threshold calibrations to
account for what the field-deformation account explains with a single
principle: each region is selectively sensitive to one operation in the
decomposition $V = \chi \circ D_\ell \circ C_{\sigma}$, and the
selectivity is structural rather than incidental.

The anterior insula is not claimed to be intrinsically a logic area. The
claim is narrower: in tasks where syntactic construction and comparison
are controlled by design, the residual negation contrast localizes to
this region, and the localization is consistent with the field-deformation
account and inconsistent with the generic-load account as formulated.
\end{proof}

\subsection{The Four-Stage Model}

The dissociation argument supports a functional architecture in which the
verification process decomposes into four stages. We state this
architecture explicitly and note its relationship to both the
field-theoretic framework and the neural evidence.

\begin{enumerate}
\item \emph{Parse}: The sentence is incrementally parsed and its
syntactic structure is recovered. This stage corresponds to $C_{\sigma}$
applied to the surface string. It is associated with left inferior
frontal and posterior temporal regions and is sensitive to syntactic
complexity but not to logical form per se.

\item \emph{Construct field}: The parsed structure is used to construct
an inferential field $\mathcal{F} = (X, R, \omega, A)$ for the sentence.
This is the stage at which the distinction space, reachability relation,
orientation field, and admissibility structure are established. It is
part of $C_{\sigma}$ in the formal decomposition but is conceptually
distinct from surface parsing: the same surface structure can, in
principle, correspond to different inferential fields depending on
lexical semantic content.

\item \emph{Deform field}: Logical operators, including negation,
quantifiers, and monotonicity-reversing expressions, transform the
constructed field. For negation, this stage applies $D_N$ to the
orientation field, reversing $\omega$ throughout the relevant domain.
This is the stage associated with $R_D$ in the neural evidence and with
the anterior insula cluster. Its cost is proportional to the orientation
distortion $\|\Delta\|$ and the boundary curvature $\kappa_\omega$.

\item \emph{Verify scenario}: The deformed field is compared against the
visual or described scenario. The verification function $\chi$ computes
whether the scenario is reachable from the sentence representation under
the current field geometry. This stage is associated with $R_\chi$ and
parietal comparison regions. Its cost includes the residual distortion
cost $\|\Delta\|_\tau$ carried over from the deformation stage, which is
why delayed verification shows sustained negation effects even when the
deformation stage has long concluded.
\end{enumerate}

\begin{corollary}[Empirical constraint on the four-stage model]
\label{cor:four-stage}
The four-stage model is empirically constrained by the triple dissociation
in Proposition \ref{prop:dissociation}. It is not a free theoretical
posit but a model whose stages correspond to functionally and at least
partially neurally separable operations. The model remains a theoretical
abstraction rather than a direct anatomical map, but it is constrained
by the anatomy in the sense that any reduction of the four stages to
fewer stages would require the elimination of one of the three
dissociable contrasts.
\end{corollary}

\begin{proof}
If stages 1 and 2 were merged into a single undifferentiated parsing
operation, then the syntactic complexity contrast and the field
construction contrast would be predicted to co-localize, with no room
for a distinct field-construction signature. If stages 3 and 4 were
merged, then the negation contrast and the comparison contrast would
co-localize, contrary to the observed anatomical separation. If stages
2 and 3 were merged, then the boundary between syntactic construction
and logical deformation would be eliminated, predicting that negation
effects and syntactic complexity effects should modulate the same
neural populations. The triple dissociation rules out all three
reductions.
\end{proof}

\subsection{The Model as Computational Interpretation}

The four-stage model is a computational interpretation of the inferential
field framework, not an additional hypothesis added to it. The framework
defines the objects --- fields, orientations, deformations, admissibility
structures --- and the model describes how processing moves through those
objects in sequence. Each stage is defined by its input, its output, and
the transformation it applies.

The model makes the processing architecture explicit in a way that
permits contact with both the behavioral and the neural literature. On
the behavioral side, it predicts that manipulations targeting different
stages should have separable effects on reaction time distributions,
error patterns, and the time course of verification difficulty. On the
neural side, it predicts that each stage should have a distinct neural
profile, with cross-stage interactions occurring at the transitions
between stages rather than uniformly throughout processing.

The most important transition for the present paper is between stages 3
and 4: between field deformation and scenario verification. The residual
distortion $\|\Delta\|$ is produced at stage 3 and consumed at stage 4.
The delayed-verification paradigm manipulates the interval between these
stages. The persistence result established in Section 5 --- that
$\|\Delta\|_\tau$ does not diminish with delay in the absence of active
re-orientation --- is therefore a claim about what happens to the output
of stage 3 during the interval before stage 4 begins. The distorted
representation is maintained in working memory without spontaneous
re-normalization. Stage 4 receives it in its distorted form and incurs
the verification cost accordingly.

This account of the stage-3 to stage-4 transition is the field-theoretic
explanation of the most theoretically challenging result in the empirical
record: the sustained cost of negation across temporal delays. It
requires no special mechanism beyond the geometry of the inferential
field and the claim that working memory preserves rather than normalizes
the structural properties of stored representations.

\section{Connections to Admissibility Geometry and Distinguishability
Theory}

The inferential field framework developed in this paper did not emerge in
isolation. It belongs to a broader research program organized around a
cluster of related formal structures: admissibility geometry, which
formalizes the conditions under which a state or inference is available
for continuation; distinguishability theory, which formalizes the
conditions under which a difference between states can be represented and
acted upon; and repair theory, which formalizes the conditions under
which a damaged or incomplete representation can be restored to
admissible form. The present section makes the connections between these
programs explicit. The goal is not to subsume the inferential field
account into a larger system by stipulation but to show that the
mathematical structure developed to handle negation and monotonicity is
the same structure that appears independently in memory, repair, and
scientific explanation. This convergence is evidence that the structure
is tracking something real.

\subsection{The Admissibility Program}

The admissibility program begins from the observation that the standard
epistemological and semantic frameworks begin too late. They ask which
beliefs are true, which inferences are valid, which representations are
accurate. But prior to any of these questions is a more fundamental one:
which states are available for cognitive engagement at all? A belief that
cannot be entertained, an inference that cannot be drawn, a
representation that cannot be accessed are not false or invalid; they
are inadmissible. The admissibility structure $A$ of the inferential
field is the local instantiation of this more general concept.

In the admissibility program, the central formal object is the
\emph{admissibility manifold}: the space of states that are locally
available for continuation, equipped with a geometry that reflects how
admissibility conditions vary across the state space. A state is
admissible if it lies within this manifold; it becomes inadmissible when
it is deformed outside the manifold by some transformation. Repair is
the process of returning an inadmissible state to the manifold.

The inferential field connects to this program through the admissibility
structure $A : X \to \{0,1\}$. An admissible region in the inferential
field --- a connected component of $A^{-1}(1)$ --- is a local patch of
the admissibility manifold restricted to the distinction space $X$. The
licensing domain for NPIs defined in Section 4 is such a patch: a region
where both admissibility and orientation conditions are satisfied, and
within which NPI licensing proceeds without restriction.

The connection runs deeper than this local correspondence, however. The
admissibility program studies how admissibility conditions change under
transformation: which transformations preserve admissibility, which
deform it, and which destroy it entirely. Negation, in the inferential
field framework, is a transformation of the orientation component of the
field. It does not directly modify the admissibility structure $A$; it
reverses $\omega$ while leaving $A$ unchanged. This means that negation
can take an inferential system from an admissible orientation into an
admissible reversed orientation, or from an admissible region into an
inadmissible one if the admissibility conditions are themselves
orientation-sensitive.

This last case is precisely NPI licensing failure. When a negative
operator is embedded in a context that makes the resulting field
inadmissible --- when $N(\omega) < 0$ in a region where $A = 0$ ---
the NPI probe returns zero not because the orientation is wrong but
because the admissibility structure excludes the relevant region. The
field has been deformed into an inadmissible configuration. The grammar
registers this as ungrammaticality; the admissibility program registers
it as a deformation outside the manifold.

\begin{proposition}[Negation as admissibility-preserving transformation]
\label{prop:adm-preserving}
The negation operator $N : \mathcal{F} \to \mathcal{F}$ preserves the
admissibility structure: $N(A) = A$. Negation can nevertheless produce
admissibility failure when the admissibility structure is
orientation-sensitive, that is, when $A(x)$ depends on $\omega(x)$.
\end{proposition}

\begin{proof}
By definition, $N(\omega)(x) = -\omega(x)$ and $N$ leaves $A$
unchanged. So $N(A) = A$ directly. However, if $A$ is defined in terms
of orientation --- if, for instance, $A(x) = 1$ only when $\omega(x)
< 0$ --- then after applying $N$, the effective admissibility of a
region that was previously admissible may change, because $\omega$ has
been reversed. The admissibility function $A$ as a formal object is
unchanged; but the conditions that must be met for $A(x) = 1$ to obtain
may now fail in regions that were previously admissible. This is the
formal counterpart of the empirical observation that negation can both
create and destroy NPI licensing environments: it creates downward
orientation where there was none, and destroys it where it previously
existed, while the structural admissibility conditions remain fixed.
\end{proof}

\subsection{Distinguishability Geometry}

Distinguishability geometry begins from a complementary starting point.
Where the admissibility program asks which states are available for
continuation, distinguishability geometry asks which differences between
states can be represented. A distinction that cannot be represented is
not false; it is indistinguishable. The formal object is a
\emph{distinguishability space} $(X, \sim)$, where $\sim$ is an
equivalence relation encoding which pairs of states are
indistinguishable under the current representational resources.

The inferential field connects to distinguishability geometry through the
distinction space $X$ itself. The elements of $X$ are distinctions: they
are, by definition, differences that the system can represent. The
reachability relation $R$ on $X$ encodes which distinctions are
inferentially connected: which differences can be reached from which
others by inference. A transformation that alters $R$ changes the
inferential connectivity of the distinction space; a transformation that
collapses distinctions --- merging elements of $X$ --- reduces
distinguishability directly.

Negation, in the distinguishability-theoretic perspective, is a
transformation that preserves the elements of $X$ while reversing the
orientation of their inferential connections. It does not collapse
distinctions; it reorients them. The distinctions that were reachable
in the upward direction become reachable in the downward direction, and
vice versa. The negation distortion $D_N(x) = d_R(x, N(x))$ measures
how far each distinction is displaced in the reachability structure by
this reorientation: a distinction that was easily reachable from its
neighbors may become distant after orientation reversal, if the
reachability paths that connected it depended on the upward-expansion
direction.

\begin{definition}[Distinguishability-theoretic negation distortion]
In the distinguishability-theoretic interpretation, the negation
distortion $D_N(x)$ measures the change in inferential distance between
$x$ and its inferential neighbors induced by the orientation reversal
$N$. Formally,
\[
D_N(x) = \sum_{y \in \mathcal{N}(x)} \left| d_R(x, y) - d_R(N(x),
N(y)) \right|,
\]
where $\mathcal{N}(x)$ is the set of inferential neighbors of $x$ under
$R$.
\]
\end{definition}
This definition makes explicit that the distortion is not a global
property of the negation operator but a local property of each
distinction and its neighborhood. Distinctions that are well-connected
in both orientations --- whose inferential neighbors are reachable
regardless of field orientation --- have low distortion. Distinctions
that are reachable primarily through upward-expanding paths have high
distortion under negation, because the paths that connected them are no
longer available in the reversed field.

The three cost components identified in Section 5 now have a unified
distinguishability-theoretic interpretation:

\begin{enumerate}
\item \emph{Orientation cost} is the cost of representing a distinction
whose inferential neighborhood has been reoriented: the system must
navigate connections in the non-default direction.
\item \emph{Boundary cost} is the cost of representing a distinction
near a distinguishability boundary: a region where the inferential
distance between neighboring distinctions changes rapidly as orientation
varies.
\item \emph{Residual distortion cost} is the cost of maintaining a
representation whose inferential neighborhood is displaced from its
default configuration: the working-memory system must actively preserve
an anomalous reachability structure against normalization.
\end{enumerate}

\subsection{Repair Theory}

Repair theory addresses the conditions under which an inadmissible or
damaged representation can be restored to admissible form. The central
formal object is the \emph{repair pseudometric} $d_{\mathrm{rep}}(x, y)$,
which measures the cost of restoring a representation $x$ to an
admissible state $y$. Repair is not mere correction; it is the process
of navigating from an inadmissible region of the state space back to the
admissibility manifold along a path of minimal cost.

The connection between repair theory and the inferential field framework
runs through the verification stage of the four-stage model. When a
negated representation $\mathcal{F}^-$ is compared against a positive
scenario $P$, the verification process must either re-orient the field
--- applying $N$ to convert $\mathcal{F}^-$ back to $\mathcal{F}^+$ ---
or navigate the comparison in the non-default direction. The first
strategy is a repair operation: it restores the representation to the
orientation that the verification process is designed to handle. The cost
of this repair is proportional to $\|\Delta\|$, the global negation
distortion. The repair pseudometric gives a formal measure of this cost:
\[
d_{\mathrm{rep}}(\mathcal{F}^-, \mathcal{F}^+) = \|\Delta\| =
\int_X D_N(x)\, d\mu(x).
\]
This identification has a theoretical consequence. In repair theory, the
repair pseudometric satisfies a conservation law: the total repair cost
of a sequence of transformations is bounded below by the distortion
accumulated along the sequence. Applied to the verification process, this
means that the total cost of verifying a negated sentence against a
scenario is bounded below by $\|\Delta\|$ regardless of verification
strategy. The bound cannot be avoided by choosing a more efficient
verification algorithm; it is a consequence of the geometry of the
inferential field.

\begin{proposition}[Lower bound on verification cost]
\label{prop:lower-bound}
For any verification strategy $\chi$ and any negated representation
$\mathcal{F}^-$,
\[
C_V(\mathcal{F}^-, P) \geq C_V(\mathcal{F}^+, P) +
d_{\mathrm{rep}}(\mathcal{F}^-, \mathcal{F}^+),
\]
where $C_V(\mathcal{F}, P)$ is the cost of verifying scenario $P$
against field $\mathcal{F}$.
\end{proposition}
\begin{proof}
The verification cost for $\mathcal{F}^-$ includes at minimum the cost
of verifying $\mathcal{F}^+$ plus the cost of resolving the distortion
between $\mathcal{F}^-$ and $\mathcal{F}^+$. The distortion resolution
cost is bounded below by $d_{\mathrm{rep}}(\mathcal{F}^-, \mathcal{F}^+)$
by the definition of the repair pseudometric as the minimum cost path
from $\mathcal{F}^-$ to $\mathcal{F}^+$ in the state space. No
verification strategy can avoid incurring this cost, since any strategy
that correctly computes the verification result must at some point resolve
the orientation difference between the negated representation and the
default-oriented scenario representation. The triangle inequality on
$d_{\mathrm{rep}}$ then gives the stated bound.
\end{proof}

\subsection{The RSVP Connection}

The RSVP framework --- Relativistic Scalar-Vector Plenum --- is a
field-theoretic approach to cognitive and physical dynamics that
identifies cognitive states with configurations of a scalar field $\Phi$
and a vector field $\vec{v}$ subject to constraint surfaces $S$. The
connections between RSVP and the inferential field framework are more
than analogical; they share a mathematical structure.

The correspondence is as follows:

\begin{align*}
\Phi &\longleftrightarrow \text{available inferential capacity in } X,\\
\vec{v} &\longleftrightarrow \text{inferential flow along } R,\\
S &\longleftrightarrow \text{admissibility and distinction-preserving
constraints in } A,\\
\omega &\longleftrightarrow \text{orientation of inferential flow,}\\
\kappa_\omega &\longleftrightarrow \text{curvature of constraint
surfaces.}
\end{align*}

Under this correspondence, upward entailment corresponds to a region
where inferential flow $\vec{v}$ is divergent: distinctions expand
outward as inference proceeds. Downward entailment corresponds to a
region where $\vec{v}$ is convergent: distinctions contract inward.
Negation is a local reversal of $\vec{v}$, changing divergent flow to
convergent flow in the affected region. The boundary between upward and
downward regions is a surface where $\nabla \cdot \vec{v} = 0$: neither
divergent nor convergent, the point of maximum orientation curvature
$\kappa_\omega$.

The processing cost function $C(x) = \alpha \cdot \mathbf{1}[\omega(x)
< 0] + \beta \cdot \kappa_\omega(x)$ has a natural RSVP interpretation.
The first term is the cost of operating in a convergent flow regime; the
second is the cost of operating near a constraint surface boundary where
the flow transitions between regimes. In RSVP dynamics, the energetic
cost of a trajectory through state space is determined by the integral
of the constraint violation along the path --- exactly the integral of
$C(x)$ over the inferential path developed here.

\begin{proposition}[RSVP formulation of the Orientation-Reversal
Principle]
\label{prop:rsvp}
Under the RSVP correspondence, the Orientation-Reversal Principle
(Theorem \ref{thm:orp}) states that the cognitive cost of a semantic
operator is determined by the integral of constraint violation along the
inferential trajectory induced by that operator, where constraint
violation is measured by the divergence of inferential flow and the
curvature of admissibility constraint surfaces.
\end{proposition}

\begin{proof}
By the RSVP correspondence, the orientation reversal $\int_X |\omega(x)
- T(\omega)(x)|\, d\mu(x)$ corresponds to the integral of the change in
flow divergence $|\nabla \cdot \vec{v}(x) - \nabla \cdot
T(\vec{v})(x)|$ over the affected region. The boundary curvature
$\int_X \kappa_\omega(x)\, d\mu(x)$ corresponds to the integral of
constraint surface curvature. In RSVP dynamics, the energetic cost of
a transformation is the sum of these integrals, since both divergence
change and surface curvature contribute to constraint violation along
the trajectory. This is the RSVP formulation of condition (i) and (ii)
of Theorem \ref{thm:orp}.
\end{proof}

\subsection{Convergence of Three Programs}

The convergence of the admissibility program, distinguishability
geometry, and repair theory on the same formal structure --- inferential
fields, orientation, distortion, admissibility --- is not a coincidence
of notation. It reflects the fact that all three programs are studying
the same underlying phenomenon from different starting points: the
conditions under which a cognitive system can move through a space of
distinctions while maintaining the structural properties required for
that movement to constitute reasoning.

The admissibility program asks: which states are available for
continuation? The distinguishability program asks: which differences can
be represented? The repair program asks: how can damaged representations
be restored? The inferential field provides a single object in which all
three questions can be posed simultaneously: the admissibility structure
$A$ answers the first, the distinction space $X$ and reachability
relation $R$ answer the second, and the repair pseudometric on the space
of fields answers the third. Negation, in this unified framework, is not
a logical operator that happens to have interesting cognitive
consequences. It is a transformation of the inferential field that
simultaneously affects admissibility, distinguishability, and
repairability --- and whose cognitive cost is a measure of all three.

The deeper philosophical claim, which we develop in the final section,
is that this unified structure reveals something about the relationship
between negation and logic more generally. Logic does not begin with
truth values and add negation as one operator among others. It begins
with the organization of distinctions, and negation is the fundamental
transformation that reveals the orientation of that organization. Truth,
on this view, is a derived notion: it is what you get when you have
navigated an inferential field successfully and arrived at a verified
distinction. Negation is prior to truth because the orientation structure
of the field is prior to verification.

\section{Negation Before Logic}

The preceding sections have developed a geometric account of negation
processing and connected it to three independent formal programs. The
final section draws out the philosophical consequences of this account.
The central claim is not merely that negation has an interesting
geometry, or that processing cost can be decomposed into separable
components, or even that the inferential field unifies four bodies of
experimental evidence. The central claim is stronger: negation is prior
to logic, in the sense that the orientation structure of the inferential
field is prior to the truth-conditional content that logic operates on.
This is not a paradox. It is a consequence of taking the geometry
seriously.

\subsection{The Standard Picture and Its Inversion}

The standard picture of the relationship between negation and logic runs
as follows. Logic begins with propositions, which have truth values.
Negation is an operator that maps truth values to their complements:
true becomes false, false becomes true. The cognitive cost of negation
is then the cost of computing this mapping and tracking its effects
through a complex sentence. Processing difficulty arises because the
mapping must be applied and its effects must be propagated through the
semantic composition.

This picture is coherent, but it places truth at the foundation and
treats negation as a derived operation defined over it. The empirical
evidence reviewed in this paper suggests that this ordering is wrong,
or at least incomplete. Negation effects are not confined to the
verification stage, where truth values are computed against scenarios.
They arise during field construction and deformation, before any
truth-conditional comparison has taken place. The cost of negation is
incurred in the structure of the representation, not merely in its
evaluation. This means that negation is doing something to the
inferential field that is independent of, and prior to, the computation
of truth values.

The inversion we propose is this: logic does not begin with truth and
add negation. It begins with the organization of distinctions and their
inferential connections. Negation is the transformation that reveals the
orientation of that organization. Truth is a derived notion: it is the
result of a successful navigation of the inferential field, a verified
reachability relation between a representation and a scenario. Negation
is prior to truth because orientation is prior to verification.

\subsection{Distinctions Before Propositions}

The foundation of the inferential field is the distinction space $X$.
The elements of $X$ are not propositions. They are differences: the
difference between red and not-red, between all and some, between inside
and outside, between reachable and unreachable. Propositions are
constructed from distinctions by the process of field construction
($C_\sigma$ in the four-stage model): they are organized packages of
distinctions with a specific inferential connectivity imposed by
syntactic and semantic composition.

This ordering --- distinctions before propositions --- is not merely a
formal choice. It reflects a substantive claim about what cognitive
systems represent. A system that represents distinctions is sensitive to
differences in its environment. A system that represents propositions is
sensitive to truth conditions. The claim that distinctions are prior is
the claim that sensitivity to difference is more fundamental than
sensitivity to truth. A system can represent the difference between red
and not-red without representing the proposition that something is red;
it cannot represent the proposition that something is red without
representing the distinction between red and not-red.

Negation, on this view, operates on distinctions, not on propositions.
The negation operator $N$ reverses the orientation of the inferential
field: it changes the direction in which distinctions are connected by
inference. This is a transformation of the structure of the distinction
space, not a mapping of truth values. The proposition that something is
not red is not the result of applying a truth-value complement to the
proposition that something is red. It is the result of constructing an
inferential field in which the distinction between red and not-red is
oriented in the reversed direction.

\begin{proposition}[Negation as distinction transformation]
\label{prop:neg-distinction}
The negation operator $N : \mathcal{F} \to \mathcal{F}$ is a
transformation of the distinction space prior to truth evaluation. Its
action on the orientation field $\omega$ is defined independently of any
assignment of truth values to elements of $X$.
\end{proposition}

\begin{proof}
By Definition 2.4, $N(\omega)(x) = -\omega(x)$ for all $x \in X$.
This definition makes no reference to a truth assignment $v : X \to
\{0,1\}$. The orientation reversal is a property of the inferential
field structure $(X, R, \omega, A)$, which is defined prior to and
independently of any verification comparison $\chi(P, \mathcal{F})$. A
truth value is assigned to a distinction only at the verification stage,
when the field is compared against a scenario. Negation operates at the
deformation stage, which precedes verification in the four-stage model.
Therefore negation is defined and operates independently of truth
evaluation.
\end{proof}

\subsection{Orientation as the Primitive Logical Notion}

If negation is prior to truth, then the primitive logical notion is not
truth but orientation. An inferential field has orientation before it
has truth conditions. The orientation field $\omega : X \to [-1, +1]$
encodes the direction in which inference flows through the distinction
space. Upward orientation means that inference expands: from specific to
general, from narrow to broad, from few to many. Downward orientation
means that inference contracts: from general to specific, from broad to
narrow, from many to few. This expansion and contraction is a property
of the field prior to any comparison with a scenario.

Logic, on the standard picture, begins with truth and defines validity
as truth preservation. On the present picture, logic begins with
orientation and defines inference as orientation-preserving motion
through the distinction space. Validity is not truth preservation; it
is orientation-consistent reachability. A valid inference is one that
moves through the field without violating the orientation constraints:
upward in upward regions, downward in downward regions, with costs
incurred at orientation boundaries.

This reinterpretation has consequences for classical logical laws. The
law of double negation elimination --- $\neg\neg p \vdash p$ --- is,
in the present framework, the statement that $N \circ N =
\mathrm{id}_{\mathcal{F}}$ (Proposition \ref{prop:double-neg}). This is
not an axiom added to the logic; it is a theorem of the geometry. The
law of excluded middle --- $p \vee \neg p$ --- is, in the present
framework, the statement that every distinction $x \in X$ lies either
in the upward regime ($\omega(x) > 0$) or the downward regime
($\omega(x) < 0$), with the boundary $\omega(x) = 0$ corresponding to
the classical gap. Whether the excluded middle holds in the present
framework depends on whether $\omega$ is constrained to avoid zero:
if $\omega : X \to [-1,+1] \setminus \{0\}$, then excluded middle
holds by fiat; if $\omega$ is permitted to reach zero, then there are
orientation-boundary distinctions for which neither upward nor downward
inference is available, corresponding to a failure of excluded middle.

Intuitionistic logic, which rejects excluded middle, corresponds in the
present framework to an inferential field in which $\omega$ is permitted
to take the value zero and in which constructions at the orientation
boundary are genuinely available. This is not merely an analogy; it is
a geometric interpretation of the constructive reading of logic. A
proposition is constructively provable if and only if it is reachable
through an orientation-consistent path from the available distinctions;
it is not merely the case that its negation is unreachable.

\subsection{The Priority of Admissibility}

If orientation is the primitive logical notion and truth is derived from
successful field navigation, then admissibility is prior to both. A
distinction must be admissible --- it must lie in a region where $A(x)
= 1$ --- before its orientation can be evaluated and before any
inferential motion through it can occur. Inadmissible distinctions are
not false; they are not available for inferential engagement. They lie
outside the region of the field in which logic operates.

This priority of admissibility has a consequence for the traditional
epistemological question of how we know what we know. The standard
question is: which of our beliefs are true, and how do we come to know
them? The admissibility-first picture replaces this with a prior
question: which distinctions are available for cognitive engagement, and
what determines the boundaries of availability? Truth and knowledge are
questions that arise only within the admissible region of the
inferential field. Outside that region, there are distinctions that
cannot be entertained, inferences that cannot be drawn, and
representations that cannot be accessed. The cognitive system does not
represent these distinctions as false; it does not represent them at
all.

This is the connection to the broader admissibility program: the
admissibility structure $A$ of the inferential field is the local
instantiation of a general principle that admissibility is prior to
truth. The epistemological consequences of this principle extend far
beyond the processing of negation. But the processing of negation
provides the clearest and most directly empirically constrained
illustration of the principle, because negation is the operation that
most visibly tests the boundary between the admissible and the
inadmissible. An NPI in an unlicensed environment is not merely
ungrammatical; it is a probe that has encountered an inadmissible
region and returned zero. The grammar's sensitivity to this is a
reflection of the cognitive system's sensitivity to admissibility
boundaries.

\subsection{Truth as Verified Reachability}

The preceding subsections have argued that orientation is prior to truth
and admissibility is prior to orientation. We can now state the positive
account of truth that emerges from this ordering.

\begin{definition}[Truth as verified reachability]
A distinction $x \in X$ is \emph{true in scenario} $P$ if and only if
$x$ is admissible ($A(x) = 1$), the local orientation at $x$ is
consistent with the inferential path from the sentence representation to
$x$ ($\omega$ is not violated along the path), and $x$ is reachable
from the scenario representation under the verification function $\chi$.
Formally,
\[
\mathrm{True}(x, P) \iff A(x) = 1 \;\wedge\; \chi(P, x) = 1.
\]
\end{definition}
This definition makes truth a three-stage achievement: first,
admissibility must be satisfied; second, inferential navigation to $x$
must be orientation-consistent; third, the scenario must verify $x$ as
reachable. The familiar truth conditions of propositional and predicate
logic are recovered as the special case where $A$ is everywhere 1
(everything is admissible), $\omega$ is everywhere positive (all
inference is upward-monotone), and $\chi$ is classical bivalent
comparison. The present framework generalizes this special case by
allowing $A$ to exclude regions, $\omega$ to vary continuously, and
$\chi$ to register degrees of reachability rather than binary
verification.

Negation, on this account, does not change the truth value of a
proposition by flipping a bit. It changes the orientation of the
inferential path to $x$, which changes the verification conditions:
the scenario that verifies $x$ in a downward-oriented context is
different from the scenario that verifies $x$ in an upward-oriented
context, because the reachability conditions depend on the orientation
of the field. Truth is not a static property of propositions but a
dynamic property of verification relations between oriented
representations and scenarios.

\subsection{Negation Before Logic: The Thesis}

We can now state the central philosophical thesis of this paper with
precision.

\begin{theorem}[Negation before logic]
\label{thm:nbl}
In the inferential field framework, negation is defined at the level of
orientation structure, prior to the assignment of truth values and
independent of the propositional content of the sentence. The logical
laws governing negation --- double negation elimination, the relationship
between upward and downward entailment, the licensing conditions of
negative polarity items --- are theorems of the geometry of inferential
fields, not axioms of a logical system. Truth is derived from successful
field navigation under admissibility and orientation constraints.
Admissibility is prior to truth; orientation is prior to truth;
negation, as an orientation transformation, is prior to truth.
\end{theorem}

\begin{proof}
The components of this theorem have been established in the preceding
sections. Negation is defined as $N(\omega)(x) = -\omega(x)$,
independently of truth evaluation (Proposition
\ref{prop:neg-distinction}). Double negation elimination is a theorem
of the geometry ($N \circ N = \mathrm{id}$, Proposition
\ref{prop:double-neg}). The relationship between upward and downward
entailment is a theorem of the orientation field (Proposition
\ref{prop:orientation}). NPI licensing is derived from the joint
admissibility and orientation conditions of Definition 4.1 and explained
by Proposition \ref{prop:npi-contrib}. Truth is defined as verified
reachability (Definition 8.1), which presupposes admissibility and
orientation-consistent navigation. Admissibility is prior because
$A(x)$ must equal 1 before any inferential operation on $x$ is defined.
Orientation is prior because $\omega$ is defined on the field before
$\chi$ is applied. Negation as orientation transformation is therefore
prior to truth as verified reachability.
\end{proof}

\subsection{Consequences and Open Questions}

The thesis of negation before logic has consequences that extend beyond
the scope of this paper. We close by noting the most significant open
questions that the framework generates.

The first concerns the relationship between the inferential field and
formal proof theory. If the logical laws governing negation are theorems
of the geometry rather than axioms, then there should be a precise
correspondence between the geometry of inferential fields and the proof
theory of the relevant logical system. The double-negation result
($N \circ N = \mathrm{id}$) corresponds to the eliminability of double
negation. The orientation boundary at $\omega = 0$ corresponds to the
constructive gap in intuitionistic logic. Establishing this
correspondence formally --- showing that the category of inferential
fields is equivalent to the category of models of some proof-theoretically
characterized logic --- is a substantial mathematical project that the
present paper leaves open.

The second concerns the neural basis of admissibility. The four-stage
model identifies field construction, deformation, and verification as
partially neurally separable operations. But the admissibility structure
$A$ does not map cleanly onto any of these stages; it is a constraint
that operates throughout the process. Identifying the neural correlates
of admissibility --- the mechanisms that determine which regions of the
distinction space are open for inferential engagement at any given moment
--- is an open empirical question. The NPI licensing literature provides
indirect evidence, since NPI licensing failure is a behavioral signature
of inadmissibility, but direct neural evidence for an admissibility
mechanism is currently lacking.

The third concerns the extension to non-linguistic inference. The
inferential field framework has been developed here in the context of
natural language processing, where the empirical record is richest. But
the formal structure --- distinctions, reachability, orientation,
admissibility --- is not intrinsically linguistic. If the same structure
appears in visual reasoning, spatial navigation, conceptual combination,
or scientific inference, then the framework has implications beyond
language. The admissibility and distinguishability programs have already
extended similar formal structures to memory and repair; the extension
to non-linguistic inference would complete the unification.

The fourth concerns the development of experimental predictions that
distinguish the field-theoretic account from its competitors in cases
not yet tested. The most promising direction is the construction of
sentences that traverse orientation domain boundaries without using
negative morphemes: purely quantificational sentences whose monotonicity
profile includes a boundary between upward and downward regions. The
field-theoretic account predicts that such sentences should incur
boundary costs even in the absence of any negative expression. This
prediction is not derivable from operator-count or polarity accounts and
would, if confirmed, provide direct evidence for the curvature term
$\kappa_\omega$ as an independent contributor to processing cost.

The inferential field is not the completion of a theory. It is the
beginning of one. The phenomena that motivated it --- the cost of
negation, the geometry of monotonicity, the licensing of polarity items,
the persistence of orientation in working memory --- are the clearest
cases. The framework predicts that the same geometry will appear wherever
inference is constrained by orientation: in the resolution of
presuppositions, in the processing of conditionals, in the updating of
beliefs under contradictory evidence. Whether it does is a question that
the framework has now made precise enough to answer empirically.


\newpage 
\section*{Appendices}

\appendix

\section{Monotonicity as Reachability Preservation}

\subsection{Reachability Spaces}

\begin{definition}[Reachability Space]
A \emph{reachability space} is a pair
\[
\mathcal{R}=(X,R)
\]
where X is a set and
\[
R\subseteq X\times X
\]
is a preorder satisfying
\begin{enumerate}
\item Reflexivity:
\[
R(x,x)
\]
for all x\in X.
\item Transitivity:
\[
R(x,y)\wedge R(y,z)
\Rightarrow
R(x,z).
\]
\end{enumerate}
\end{definition}
For x\in X, define the forward reachability set

\[
\mathcal{N}^{+}(x)
{y\in X:R(x,y)}.
\]

Similarly define

\[
\mathcal{N}^{-}(x)
{y\in X:R(y,x)}.
\]

\begin{definition}[Inferential Inclusion]
For subsets A,B\subseteq X,

\[
A\preceq_R B
\]
iff

\[
A\subseteq B.
\]
The induced expansion operator is

\[
\mathcal{E}(A)
{y:\exists x\in A,;R(x,y)}.
\]
\end{definition}

\subsection{Monotonicity Transformations}

\begin{definition}[Orientation Preserving Map]
A map

\[
T:X\rightarrow X
\]
is orientation preserving iff

\[
R(x,y)
\Longrightarrow
R(T(x),T(y)).
\]
\end{definition}
\begin{definition}[Orientation Reversing Map]
A map

\[
T:X\rightarrow X
\]
is orientation reversing iff

\[
R(x,y)
\Longrightarrow
R(T(y),T(x)).
\]
\end{definition}
\begin{definition}[Monotone Operator]
A transformation T is monotone iff it preserves inferential inclusion:

\[
A\subseteq B
\Longrightarrow
T(A)\subseteq T(B).
\]
\end{definition}
\begin{definition}[Antitone Operator]
A transformation T is antitone iff

\[
A\subseteq B
\Longrightarrow
T(B)\subseteq T(A).
\]
\end{definition}
\subsection{Reachability Characterization}

\begin{theorem}[Monotonicity-Reachability Equivalence]
Let T:X\to X.

Then T is monotone iff T is orientation preserving.
\end{theorem}

\begin{proof}
Assume T is orientation preserving.

Let

\[
A\subseteq B.
\]
Take

\[
y\in T(A).
\]
Then

\[
y=T(x)
\]
for some x\in A.

Since A\subseteq B,

\[
x\in B,
\]
hence

\[
y\in T(B).
\]
Therefore

\[
T(A)\subseteq T(B).
\]
Thus T is monotone.

Conversely suppose T is monotone.

Let

\[
R(x,y).
\]
Then

\[
{x}\subseteq {x,y}.
\]
Monotonicity implies

\[
T({x})
\subseteq
T({x,y}).
\]
Therefore

\[
R(T(x),T(y))
\]
and orientation is preserved.
\end{proof}

\begin{theorem}[Antitonicity-Reversal Equivalence]
A transformation is antitone iff it reverses orientation.
\end{theorem}

\begin{proof}
Suppose

\[
A\subseteq B.
\]
Antitonicity gives

\[
T(B)\subseteq T(A).
\]
Applying this to singleton-generated neighborhoods yields

\[
R(x,y)
\Longrightarrow
R(T(y),T(x)).
\]
The converse follows identically.
\end{proof}

\subsection{Negation as Orientation Reversal}

\begin{definition}[Negation Transformation]
A negation operator is an involutive antitone map

\[
N:X\rightarrow X
\]
satisfying

\[
N^2=I.
\]
\end{definition}
\begin{lemma}
For every reachability relation R,

\[
R(x,y)
\Longrightarrow
R(N(y),N(x)).
\]
\end{lemma}
\begin{proof}
Immediate from antitonicity.
\end{proof}

\begin{theorem}[Double-Reversal Restoration]
\label{app:double-restoration}
For every reachability space (X,R),

\[
N^2=I
\]
implies restoration of inferential orientation.
\end{theorem}

\begin{proof}
Let

\[
R(x,y).
\]
Applying N,

\[
R(N(y),N(x)).
\]
Applying N again,

\[
R(N^2(x),N^2(y)).
\]
Since

\[
N^2=I,
\]
we obtain

\[
R(x,y).
\]
Thus original orientation is recovered.
\end{proof}

\begin{corollary}[Double Negation Preservation]
For every subset A\subseteq X,

\[
N(N(A))
A.
\]
\end{corollary}

\begin{proof}
Directly from N^2=I.
\end{proof}

\subsection{Orientation Classes}

Define

\[
\mathcal{M}^{+}
{T:T \text{ preserves orientation}}
\]

and

\[
\mathcal{M}^{-}
{T:T \text{ reverses orientation}}.
\]

\begin{proposition}
\[
\mathcal{M}^{+}
\]
forms a monoid under composition.
\end{proposition}

\begin{proof}
Identity preserves orientation.

Composition of orientation-preserving maps preserves orientation.

Associativity follows from function composition.
\end{proof}

\begin{proposition}
The composition law satisfies

\[
(+)(+)=(+),
\]
\[
(+)(-) = (-),
\]
\[
(-)(+) = (-),
\]
\[
(-)(-) = (+).
\]
\end{proposition}
\begin{proof}
Immediate from repeated application of orientation-preservation and orientation-reversal definitions.
\end{proof}

\begin{corollary}[Parity Law]
The net orientation of a sequence of negations depends only on the parity of the number of reversals.
\end{corollary}

\begin{proof}
Repeated application of

\[
(-)(-)=(+).
\]
Even parity yields preservation.

Odd parity yields reversal.
\end{proof}

\subsection{Reachability Interpretation of Monotonicity}

Classical monotonicity theory is therefore recovered as a special case of orientation behavior on reachability spaces:

\[
\text{upward-entailing}
\iff
\text{orientation preserving},
\]
\[
\text{downward-entailing}
\iff
\text{orientation reversing}.
\]
Negation is not fundamentally a truth-value operator.

It is an involutive reversal of inferential orientation.

Monotonicity is therefore a property of reachability geometry rather than lexical polarity.

\section{Geometry of Orientation Boundaries}

\subsection{Orientation Manifolds}

Let

\[
(M,g)
\]
be a smooth connected Riemannian manifold.

\begin{definition}[Orientation Field]
An orientation field is a smooth scalar function

\[
\omega:M\rightarrow\mathbb{R}.
\]
Positive values correspond to orientation-preserving regions, negative values correspond to orientation-reversing regions, and the boundary set is

\[
\partial\Omega
{x\in M:\omega(x)=0}.
\]
\end{definition}

Define

\[
\Omega^{+}
{x:\omega(x)>0}
\]

and

\[
\Omega^{-}
{x:\omega(x)<0}.
\]

Then

\[
M=\Omega^{+}\cup\partial\Omega\cup\Omega^{-}.
\]
\subsection{Orientation Gradient}

\begin{definition}[Orientation Gradient]
The local orientation gradient is

\[
\nabla\omega.
\]
Its magnitude

\[
G(x)
|\nabla\omega(x)|
\]

measures the rate of orientation change.
\end{definition}

\begin{definition}[Orientation Energy]
The orientation energy of a region U\subseteq M is

\[
E_\omega(U)
\int_U
|\nabla\omega|^2,dV.
\]
\end{definition}

\begin{theorem}[Minimum-Energy Orientation]
Among all fields satisfying fixed boundary conditions,

\[
E_\omega
\]
is minimized iff

\[
\Delta\omega=0.
\]
\end{theorem}
\begin{proof}
The Euler-Lagrange equation for

\[
E_\omega
\int
|\nabla\omega|^2,dV
\]

is

\[
\Delta\omega=0.
\]
Hence minimizers are harmonic.
\end{proof}

\subsection{Boundary Curvature}

Assume

\[
\nabla\omega\neq0
\]
on

\[
\partial\Omega.
\]
Then \partial\Omega is a smooth hypersurface.

\begin{definition}[Unit Orientation Normal]
\[
n
\frac{\nabla\omega}
{|\nabla\omega|}.
\]
\end{definition}

\begin{definition}[Orientation Curvature]
The orientation curvature is

\[
\kappa_\omega
\nabla\cdot n.
\]
\end{definition}

Equivalently,

\[
\kappa_\omega
\nabla\cdot
\left(
\frac{\nabla\omega}
{|\nabla\omega|}
\right).
\]

\begin{proposition}
If \partial\Omega is locally planar,

\[
\kappa_\omega=0.
\]
\end{proposition}
\begin{proof}
For a planar boundary the unit normal is constant.

Hence

\[
\nabla\cdot n=0.
\]
\end{proof}
\subsection{Curvature Concentration}

Define

\[
C(U)
\int_U
|\kappa_\omega|
,dV.
\]

\begin{definition}[Boundary Complexity]
The total boundary complexity is

\[
C_\omega
\int_{\partial\Omega}
|\kappa_\omega|
,dS.
\]
\end{definition}

\begin{theorem}
For every compact boundary,

\[
C_\omega
\geq 0.
\]
Equality holds iff every connected component of
\partial\Omega is flat.
\end{theorem}

\begin{proof}
Absolute values are nonnegative.

Vanishing integral implies

\[
|\kappa_\omega|=0
\]
almost everywhere.

Thus

\[
\kappa_\omega=0.
\]
\end{proof}
\subsection{Orientation Distortion Functional}

Let

\[
T:M\rightarrow M
\]
be a deformation.

\begin{definition}[Orientation Distortion]
\[
D_\omega(T)
\int_M
|\omega(x)-\omega(Tx)|
,dV.
\]
\end{definition}

\begin{definition}[Gradient Distortion]
\[
D_{\nabla}(T)
\int_M
|
\nabla\omega(x)

\nabla\omega(Tx)
|
,dV.
\]
\end{definition}

\begin{definition}[Curvature Distortion]
\[
D_\kappa(T)
\int_M
|\kappa_\omega(x)-\kappa_\omega(Tx)|
,dV.
\]
\end{definition}

\subsection{Boundary Cost Functional}

Define

\[
\mathcal C(T)
\alpha D_\omega(T)
+
\beta D_{\nabla}(T)
+
\gamma D_\kappa(T)
\]

for positive constants

\[
\alpha,\beta,\gamma>0.
\]
\begin{theorem}[Boundary Cost Decomposition]
Every smooth orientation deformation admits the decomposition

\[
\mathcal C(T)
\mathcal C_{\text{field}}
+
\mathcal C_{\text{gradient}}
+
\mathcal C_{\text{curvature}}.
\]
\end{theorem}

\begin{proof}
Immediate from linearity of integration.
\end{proof}

\subsection{Boundary Amplification Theorem}

\begin{theorem}
Let

\[
T_\varepsilon
I+\varepsilon V
\]

be a smooth perturbation.

Then

\[
D_\omega(T_\varepsilon)
O(\varepsilon)
\]

while

\[
D_\kappa(T_\varepsilon)
O(\varepsilon |\nabla\kappa_\omega|).
\]

Consequently regions of high curvature produce disproportionately large distortion.
\end{theorem}

\begin{proof}
Taylor expansion gives

\[
\omega(T_\varepsilon x)
\omega(x)
+
\varepsilon
\nabla\omega\cdot V
+
O(\varepsilon^2).
\]

Hence

\[
D_\omega
O(\varepsilon).
\]

Similarly,

\[
\kappa_\omega(T_\varepsilon x)
\kappa_\omega(x)
+
\varepsilon
\nabla\kappa_\omega\cdot V
+
O(\varepsilon^2).
\]

Therefore

\[
D_\kappa
O(\varepsilon |\nabla\kappa_\omega|).
\]

Regions with rapidly varying curvature amplify distortion.
\end{proof}

\subsection{Singular Boundaries}

\begin{definition}[Orientation Singularity]
A point

\[
x_0
\]
is singular iff

\[
\omega(x_0)=0
\]
and

\[
\nabla\omega(x_0)=0.
\]
\end{definition}
\begin{proposition}
Orientation curvature becomes undefined at singular points.
\end{proposition}

\begin{proof}
The normal field

\[
n
\frac{\nabla\omega}
{|\nabla\omega|}
\]

cannot be defined.
\end{proof}

\begin{definition}[Curvature Defect]
For a neighborhood \(U),

\[
\delta_\kappa(U)
\int_U
|\nabla\kappa_\omega|^2
,dV.
\]
\end{definition}

\subsection{Global Orientation Geometry}

\begin{definition}[Total Orientation Action]
\[
S_\omega
\int_M
\left(
\lambda_1
|\nabla\omega|^2
+
\lambda_2
\kappa_\omega^2
\right)
dV.
\]
\end{definition}

\begin{theorem}
Critical points of

\[
S_\omega
\]
satisfy a fourth-order geometric field equation

\[
\Delta\omega
\lambda
\Delta\kappa_\omega

0
\]

for suitable \lambda.
\end{theorem}

\begin{proof}
Applying variational calculus to

\[
S_\omega
\]
produces the Euler-Lagrange equations involving both first-order gradient terms and second-order curvature terms.

Combining them yields the stated fourth-order equation.
\end{proof}

\subsection{Orientation Boundary Principle}

Let

\[
\mathcal P(x)
\alpha |\omega(x)|
+
\beta |\nabla\omega(x)|
+
\gamma |\kappa_\omega(x)|.
\]

\begin{theorem}[Boundary Cost Principle]
Local orientation complexity is controlled by the triple

\[
(\omega,\nabla\omega,\kappa_\omega).
\]
In particular,

\[
\mathcal P(x)
\rightarrow \infty
\]
whenever gradient concentration or curvature concentration diverges.

Thus orientation boundaries act as geometric singularities whose distortion cost is determined by both gradient magnitude and boundary curvature.
\end{theorem}

\begin{proof}
Immediate from the definition of

\[
\mathcal P.
\]
Divergence of either term forces divergence of the total complexity measure.
\end{proof}

\section{Negation Distortion and Reachability Metrics}

\subsection{Metric Reachability Spaces}

Let

\[
(X,R)
\]
be a reachability space.

\begin{definition}[Reachability Path]
A reachability path from \(x) to \(y) is a finite sequence

\[
\gamma=(x_0,x_1,\ldots,x_n)
\]
such that

\[
x_0=x,
\qquad
x_n=y,
\]
and

\[
R(x_i,x_{i+1})
\]
for all \(i\).
\end{definition}

\begin{definition}[Path Length]
Let

\[
w:X\times X\rightarrow \mathbb R_{\ge0}
\]
be an edge-weight function.

The length of a path is

\[
L(\gamma)
\sum_{i=0}^{n-1}
w(x_i,x_{i+1}).
\]
\end{definition}

\begin{definition}[Reachability Metric]
The reachability metric is

\[
d_R(x,y)
\inf_{\gamma:x\leadsto y}
L(\gamma).
\]
\end{definition}

\begin{theorem}
If

\[
w(x,y)>0
\]
for all distinct reachable pairs, then

\[
d_R
\]
is a metric on each connected component.
\end{theorem}

\begin{proof}
Non-negativity follows from positivity of weights.

Identity follows because only the empty path has length zero.

Symmetry is obtained by considering the undirected closure.

Triangle inequality follows from path concatenation.
\end{proof}

\subsection{Inferential Geodesics}

\begin{definition}[Geodesic Reachability Path]
A path

\[
\gamma^\star
\]
is geodesic iff

\[
L(\gamma^\star)
d_R(x,y).
\]
\end{definition}

\begin{definition}[Inferential Geodesic Set]
\[
\Gamma(x,y)
{
\gamma:
L(\gamma)=d_R(x,y)
}.
\]
\end{definition}

\begin{proposition}
If X is finite then

\[
\Gamma(x,y)\neq\emptyset.
\]
\end{proposition}
\begin{proof}
The infimum is attained because only finitely many path lengths exist.
\end{proof}

\subsection{Negation Distortion}

Let

\[
N:X\rightarrow X
\]
be an involutive orientation reversal.

\[
N^2=I.
\]
\begin{definition}[Pointwise Negation Distortion]
\[
D_N(x)
d_R(x,N(x)).
\]
\end{definition}

\begin{definition}[Average Negation Distortion]
For probability measure \(\mu),

\[
\bar D_N
\int_X
D_N(x)
,d\mu(x).
\]
\end{definition}

\begin{definition}[Maximum Negation Distortion]
\[
D_N^{\max}
\sup_{x\in X}
D_N(x).
\]
\end{definition}

\subsection{Distortion Geometry}

\begin{theorem}[Zero Distortion Criterion]
\[
D_N(x)=0
\]
iff

\[
N(x)=x.
\]
\end{theorem}
\begin{proof}
Since \(d_R) is a metric,

\[
d_R(x,y)=0
\]
iff

\[
x=y.
\]
Substitute

\[
y=N(x).
\]
\end{proof}
\begin{definition}[Negation Orbit]
The orbit of x under negation is

\[
\mathcal O_N(x)
{x,N(x)}.
\]
\end{definition}

\begin{proposition}
Every orbit has cardinality

\[
1
\quad\text{or}\quad
2.
\]
\end{proposition}
\begin{proof}
Since

\[
N^2=I,
\]
higher-order cycles cannot occur.
\end{proof}

\subsection{Distortion Tensor}

Assume \(X) is a smooth manifold.

Let

\[
g_{ij}
\]
denote the reachability metric tensor.

\begin{definition}[Negation Distortion Tensor]
\[
T_{ij}
g_{ij}

N^\ast(g_{ij}).
\]
\end{definition}

\begin{proposition}
\[
T_{ij}=0
\]
iff negation acts as an isometry.
\end{proposition}

\begin{proof}
The pullback metric equals the original metric precisely when N is an isometry.
\end{proof}

\subsection{Integrated Distortion Action}

\begin{definition}
The total distortion action is

\[
S_D
\int_X
D_N(x)^2
,d\mu(x).
\]
\end{definition}

\begin{theorem}
\[
S_D\ge0.
\]
Equality holds iff negation is the identity transformation.
\end{theorem}

\begin{proof}
The integrand is nonnegative.

Vanishing integral implies

\[
D_N(x)=0
\]
almost everywhere.

Hence

\[
N(x)=x.
\]
\end{proof}
\subsection{Reachability Displacement}

\begin{definition}[Displacement Field]
\[
\delta_N(x)
N(x)-x.
\]
\end{definition}

For smooth coordinates,

\[
|\delta_N(x)|
d_R(x,N(x)).
\]

\begin{definition}[Mean Squared Displacement]
\[
\Sigma_N
\int_X
|\delta_N(x)|^2
,d\mu(x).
\]
\end{definition}

\begin{proposition}
\[
\Sigma_N=S_D.
\]
\end{proposition}
\begin{proof}
By definition of displacement magnitude.
\end{proof}

\subsection{Double Negation Restoration}

\begin{theorem}[Metric Restoration Theorem]
For every point \(x\),

\[
D_{N^2}(x)=0.
\]
\end{theorem}
\begin{proof}
Since

\[
N^2(x)=x,
\]
we have

\[
D_{N^2}(x)
d_R(x,N^2(x))

d_R(x,x)

0. 

\]
\end{proof}

\begin{corollary}
\[
S_{N^2}=0.
\]
\end{corollary}
\begin{proof}
The integrand vanishes identically.
\end{proof}

\subsection{Orientation-Reversal Geodesics}

\begin{definition}
The negation geodesic associated with \(x\) is

\[
\gamma_N(x)
\operatorname{argmin}_{\gamma}
L(\gamma)
\]

subject to

\[
\gamma(0)=x,
\qquad
\gamma(1)=N(x).
\]
\end{definition}
\begin{definition}[Negation Length]
\[
\ell_N(x)
L(\gamma_N(x)).
\]
\end{definition}

\begin{proposition}
\[
\ell_N(x)
D_N(x).
\]
\end{proposition}

\begin{proof}
By geodesic optimality.
\end{proof}

\subsection{Curvature of Negation}

Let

\[
\gamma_N(t)
\]
be a smooth negation geodesic.

Define

\[
v(t)
\frac{d\gamma_N}{dt}.
\]

\begin{definition}[Negation Curvature]
\[
K_N
\left|
\frac{Dv}{dt}
\right|.
\]
\end{definition}

\begin{definition}[Integrated Curvature]
\[
\mathcal K_N
\int_{\gamma_N}
K_N,ds.
\]
\end{definition}

\begin{theorem}
Straight-line negation trajectories satisfy

\[
\mathcal K_N=0.
\]
\end{theorem}
\begin{proof}
A geodesic with constant tangent satisfies

\[
\frac{Dv}{dt}=0.
\]
\end{proof}
\subsection{Spectral Distortion}

Let

\[
\Delta_R
\]
denote the reachability Laplacian.

Suppose

\[
\Delta_R \phi_k
\lambda_k \phi_k.
\]

\begin{definition}[Spectral Negation Distortion]
\[
\Lambda_N
\sum_k
\lambda_k
|
\phi_k-N(\phi_k)
|^2.
\]
\end{definition}

\begin{theorem}
\[
\Lambda_N=0
\]
iff every Laplacian eigenmode is invariant under negation.
\end{theorem}

\begin{proof}
Every term is nonnegative.

The sum vanishes only if

\[
\phi_k=N(\phi_k)
\]
for all \(k\).
\end{proof}

\subsection{Distortion Principle}

\begin{theorem}[Negation Distortion Principle]
For every metric reachability space,

\[
D_N(x)
d_R(x,N(x))
\]

measures the inferential displacement induced by orientation reversal.

Furthermore,

\[
D_{N^2}(x)=0,
\]
so double negation is a complete restoration of reachability geometry rather than a mere cancellation of logical operators.
\end{theorem}

\begin{proof}
The first statement follows by definition.

The second follows from the Metric Restoration Theorem.
\end{proof}

\section{Persistence of Negation Distortion}

\subsection{Memory Reachability Systems}

Let

\[
(X,R,\mu)
\]
be a metric reachability space equipped with probability measure \(\mu\).

Let

\[
N:X\rightarrow X
\]
denote an involutive orientation reversal.

\[
N^2=I.
\]
\begin{definition}[Memory Reachability System]
A memory reachability system is the quintuple

\[
\mathcal M
(X,R,N,M_t,\mu)
\]

where

\[
M_t:X\rightarrow X
\]
denotes the state of memory at time t.
\]
\end{definition}

\begin{definition}[Memory Encoding]
A memory encoding operator is

\[
E:X\rightarrow \mathcal H
\]
where \(\mathcal H\) is a memory state space.
\end{definition}

\begin{definition}[Memory Retrieval]
A retrieval operator is

\[
Q:\mathcal H\rightarrow X.
\]
\end{definition}
\begin{definition}[Memory Loop]
The memory loop is

\[
\mathcal L
Q\circ E.
\]
\end{definition}

Perfect persistence corresponds to

\[
Q(E(x))
x.
\]

\subsection{Distortion Memory States}

Recall

\[
D_N(x)
d_R(x,N(x)).
\]

\begin{definition}[Stored Distortion]
The stored distortion state is

\[
\Delta_t(x)
D_N(M_t(x)).
\]
\end{definition}

\begin{definition}[Global Distortion Energy]
\[
\mathcal D_t
\int_X
\Delta_t(x)^2
,d\mu(x).
\]
\end{definition}

\subsection{Persistence Operators}

\begin{definition}[Persistence Operator]
A persistence operator is a semigroup

\[
P_t:X\rightarrow X
\]
satisfying

\[
P_0=I
\]
and

\[
P_tP_s=P_{t+s}.
\]
\end{definition}
\begin{definition}[Perfect Persistence]
Persistence is perfect iff

\[
P_t(x)=x
\]
for all \(t\).
\end{definition}

\begin{definition}[Lossless Memory]
A memory system is lossless iff

\[
M_t=P_t.
\]
\]
and

\[
d_R(M_t(x),x)=0.
\]
\end{definition}
\subsection{Distortion Conservation}

\begin{theorem}[Distortion Conservation]
Suppose memory is lossless.

Then

\[
\Delta_t(x)
\Delta_0(x)
\]

for every \(t\).
\end{theorem}

\begin{proof}
Losslessness implies

\[
M_t(x)=x.
\]
Therefore

\[
\Delta_t(x)
D_N(M_t(x))

D_N(x)

\Delta_0(x).
\]
\end{proof}

\begin{corollary}
\[
\mathcal D_t
\mathcal D_0.
\]
\end{corollary}

\begin{proof}
Integrate the previous result.
\end{proof}

\subsection{Memory Decay}

\begin{definition}[Decay Generator]
A decay process is generated by

\[
L:X\rightarrow TX
\]
through

\[
\frac{dM_t}{dt}
L(M_t).
\]
\end{definition}

\begin{definition}[Exponential Forgetting]
A memory trace obeys

\[
\frac{dm}{dt}
-\lambda m.
\]
\]

with solution

\[
m(t)
m(0)e^{-\lambda t}.
\]
\end{definition}

\begin{theorem}
Under exponential forgetting,

\[
\mathcal D_t
\mathcal D_0 e^{-2\lambda t}.
\]
\end{theorem}

\begin{proof}
\[
\Delta_t
\Delta_0 e^{-\lambda t}.
\]

Squaring and integrating gives

\[
\mathcal D_t
\mathcal D_0 e^{-2\lambda t}.
\]
\end{proof}

\subsection{Ecphory Dynamics}

\begin{definition}[Ecphory Operator]
An ecphory operator is

\[
\mathcal E_t:\mathcal H\rightarrow X.
\]
\end{definition}
\begin{definition}[Ecphory Threshold]
A threshold parameter

\[
\theta>0
\]
defines retrieval success:

\[
\mathcal E_t(h)
x
\]

iff

\[
|h-E(x)|
<
\theta.
\]
\end{definition}
\begin{definition}[Retrieval Error]
\[
\varepsilon_t(x)
d_R
\big(
x,\mathcal E_t(E(x))
\big).
\]
\end{definition}

\begin{theorem}
Perfect ecphory implies

\[
\varepsilon_t(x)=0.
\]
\end{theorem}
\begin{proof}
If retrieval reconstructs \(x\) exactly,

\[
d_R(x,x)=0.
\]
\end{proof}
\subsection{Persistence of Orientation Structure}

Let

\[
\omega_t
\]
denote the orientation field represented in memory.

\begin{definition}[Orientation Persistence]
Orientation persists iff

\[
\omega_t
\omega_0.
\]
\end{definition}

\begin{theorem}[Orientation Persistence Theorem]
If memory preserves orientation structure,

\[
\omega_t=\omega_0,
\]
then negation distortion remains invariant.
\end{theorem}

\begin{proof}
Negation distortion depends only upon

\[
\omega.
\]
If

\[
\omega_t=\omega_0,
\]
then

\[
D_N^t
D_N^0.
\]
\end{proof}

\subsection{Delayed Verification Dynamics}

Let

\[
V_t
\]
denote verification state.

\begin{definition}[Verification Functional]
\[
V_t(x)
\chi(P,M_t(x))
\]

where

\[
\chi
\]
is a consistency functional.
\end{definition}

\begin{definition}[Verification Latency]
\[
\tau
\inf
{
t:
V_t(x)=1
}.
\]
\end{definition}

\begin{theorem}
If distortion persists,

\[
\Delta_t=\Delta_0,
\]
then verification latency remains bounded below by

\[
\tau
\ge
c\Delta_0
\]
for some constant \(c>0\).
\end{theorem}

\begin{proof}
Verification requires reconstruction of orientation.

The amount of reconstruction scales with stored distortion.

Hence

\[
\tau
\propto
\Delta_0.
\]
Absorb proportionality into \(c\).
\end{proof}

\subsection{Distortion Persistence Spectrum}

Let

\[
K
\]
be the memory evolution operator.

Assume

\[
K\phi_n
\lambda_n\phi_n.
\]

\begin{definition}[Persistence Spectrum]
\[
\sigma_P
{\lambda_n}.
\]
\end{definition}

\begin{definition}[Distortion Spectrum]
\[
\Sigma_D
{
\lambda_n:
\phi_n
\text{ contributes to }
D_N
}.
\]
\end{definition}

\begin{theorem}
Modes satisfying

\[
|\lambda_n|=1
\]
produce permanent distortion persistence.
\end{theorem}

\begin{proof}
Such modes do not decay under repeated application of K.
\end{proof}

\subsection{Recoverability Geometry}

\begin{definition}[Recoverability Radius]
\[
\rho(x)
\sup
{
r:
B_r(x)
\text{ is reconstructible}
}.
\]
\end{definition}

\begin{definition}[Persistence Capacity]
\[
\Phi_P
\int_X
\rho(x)
,d\mu(x).
\]
\end{definition}

\begin{theorem}
If

\[
\Phi_P
«»

D_N^{\max},
\]

then all negation distortions are recoverable.
\end{theorem}

\begin{proof}
Every distortion lies within the recoverability radius.

Hence reconstruction exists.
\end{proof}

\subsection{Delayed Verification Theorem}

\begin{theorem}[Delayed Verification Theorem]
Suppose

\[
N
\]
introduces distortion

\[
D_N(x)>0.
\]
Assume memory preserves orientation geometry:

\[
\omega_t=\omega_0.
\]
Then

\[
\Delta_t(x)
D_N(x)
\]

for all t,

and verification latency satisfies

\[
\tau(x)
\ge
cD_N(x).
\]
Consequently delayed-verification effects arise naturally from persistence of orientation distortion rather than from repeated syntactic computation.
\end{theorem}

\begin{proof}
Distortion conservation gives

\[
\Delta_t(x)=D_N(x).
\]
The latency bound follows from the previous theorem.

Thus verification cost scales with persistent geometric distortion stored in memory.
\end{proof}

\subsection{Persistence Principle}

\begin{theorem}[Persistence Principle]
Memory preserves not propositions but distortions of reachability geometry.

Negation-induced displacement survives in memory whenever orientation structure survives.

Therefore the persistence of verification cost is equivalent to the persistence of geometric distortion.
\end{theorem}

\begin{proof}
All preceding results establish that distortion is a function of orientation geometry and that orientation persistence implies distortion persistence.

Verification cost depends on distortion.

Hence persistence of cost follows directly from persistence of geometry.
\end{proof}

\section{NPI Licensing as Admissibility Geometry}

\subsection{Admissibility Spaces}

Let

\[
(M,g)
\]
be a connected Riemannian manifold.

Let

\[
\omega:M\rightarrow\mathbb R
\]
denote an orientation field.

\begin{definition}[Admissibility Field]
An admissibility field is a smooth function

\[
A:M\rightarrow[0,1].
\]
\end{definition}
\begin{definition}[Admissible Region]
The admissible region is

\[
\mathcal A
{x\in M:A(x)>0}.
\]
\end{definition}

\begin{definition}[Fully Admissible Region]
\[
\mathcal A_1
{x:A(x)=1}.
\]
\end{definition}

\begin{definition}[Admissibility Boundary]
\[
\partial\mathcal A
{x:A(x)=0}.
\]
\end{definition}

\subsection{Licensing Domains}

\begin{definition}[Licensing Domain]
The licensing domain is

\[
L
{
x\in M:
A(x)=1,
,
\omega(x)<0
}.
\]
\end{definition}

\begin{definition}[Positive Domain]
\[
P
{
x:
A(x)=1,
,
\omega(x)>0
}.
\]
\end{definition}

\begin{definition}[Neutral Boundary]
\[
B
{
x:
A(x)=1,
,
\omega(x)=0
}.
\]
\end{definition}

Thus

\[
\mathcal A_1
P\cup B\cup L.
\]

\subsection{NPI Probe Functions}

\begin{definition}[NPI Probe]
An NPI is represented by a probe function

\[
\Pi:M\rightarrow{0,1}.
\]
\end{definition}
\begin{definition}[Licensing Condition]
A probe is licensed at x iff

\[
\Pi(x)=1.
\]
\]
and
\[
x\in L.
\]
\end{definition}
Equivalently,

\[
\Pi(x)
\chi_L(x),
\]

where

\[
\chi_L
\]
is the characteristic function of the licensing region.

\subsection{Geometric Licensing Criterion}

\begin{theorem}[Licensing Criterion]
An NPI is licensed at x iff

\[
A(x)=1
\]
and

\[
\omega(x)<0.
\]
\end{theorem}
\begin{proof}
By definition of L,

\[
x\in L
\]
iff both conditions hold.

Since

\[
\Pi=\chi_L,
\]
licensing follows immediately.
\end{proof}

\subsection{Connected Licensing Components}

\begin{definition}[Licensing Component]
A licensing component is a connected component

\[
L_i
\subseteq
L.
\]
\end{definition}
Then

\[
L
\bigsqcup_i
L_i.
\]

\begin{theorem}[Component Licensing Theorem]
Every point of a licensing component licenses NPIs.
\end{theorem}

\begin{proof}
If

\[
x\in L_i,
\]
then

\[
x\in L.
\]
Hence

\[
A(x)=1,
\qquad
\omega(x)<0.
\]
Therefore licensing holds.
\end{proof}

\subsection{Boundary Geometry}

Define

\[
\kappa_A
\nabla\cdot
\left(
\frac{\nabla A}
{|\nabla A|}
\right).
\]

\begin{definition}[Admissibility Curvature]
\[
\kappa_A
\]
is the admissibility curvature.
\end{definition}

Similarly define orientation curvature

\[
\kappa_\omega.
\]
\begin{definition}[Licensing Boundary Curvature]
\[
\kappa_L
|\kappa_A|
+
|\kappa_\omega|.
\]
\end{definition}

\begin{theorem}
Licensing uncertainty increases monotonically with

\[
\kappa_L.
\]
\end{theorem}
\begin{proof}
Large curvature implies rapid local variation of either admissibility or orientation.

Therefore infinitesimal perturbations alter membership in L.

Hence uncertainty increases.
\end{proof}

\subsection{Distance to Licensing}

\begin{definition}[Licensing Distance]
\[
d_L(x)
\inf_{y\in L}
d(x,y).
\]
\end{definition}

\begin{definition}[Licensing Potential]
\[
V_L(x)
e^{-d_L(x)}.
\]
\end{definition}

\begin{theorem}
\[
V_L(x)
1
\]

iff

\[
x\in \overline L.
\]
\end{theorem}
\begin{proof}
Distance vanishes precisely on the closure.

Hence

\[
e^{-0}=1.
\]
\end{proof}
\subsection{Licensing Flux}

Let

\[
n
\]
denote the outward normal of a licensing boundary.

\begin{definition}[Licensing Flux]
\[
\Phi_L
\int_{\partial L}
\nabla\omega\cdot n
,dS.
\]
\end{definition}

\begin{theorem}
Positive licensing flux corresponds to outward expansion of licensing domains.
\end{theorem}

\begin{proof}
By the divergence theorem,

\[
\Phi_L
\int_L
\Delta\omega
,dV.
\]

Positive flux implies net outward orientation flow.
\end{proof}

\subsection{Topological Licensing Structure}

Let

\[
H_k(L)
\]
denote singular homology groups.

\begin{definition}[Licensing Connectivity Number]
\[
\beta_0
\operatorname{rank}(H_0(L)).
\]
\end{definition}

\begin{definition}[Licensing Loop Number]
\[
\beta_1
\operatorname{rank}(H_1(L)).
\]
\end{definition}

\begin{theorem}
\[
\beta_0
\]
equals the number of disconnected licensing regions.
\end{theorem}

\begin{proof}
Standard algebraic topology.
\end{proof}

\begin{theorem}
Nonzero

\[
\beta_1
\]
implies cyclic licensing structure.
\end{theorem}

\begin{proof}
A nontrivial first homology group contains non-contractible loops.
\end{proof}

\subsection{Spectral Licensing Theory}

Let

\[
\Delta_L
\]
be the Laplace-Beltrami operator restricted to L.

Assume

\[
\Delta_L\phi_n
\lambda_n\phi_n.
\]

\begin{definition}[Licensing Spectrum]
\[
\sigma(L)
{\lambda_n}.
\]
\end{definition}

\begin{definition}[Licensing Capacity]
\[
C_L
\sum_n
\frac1{1+\lambda_n}.
\]
\end{definition}

\begin{theorem}
Larger connected licensing regions possess larger licensing capacity.
\end{theorem}

\begin{proof}
Expansion lowers low-frequency eigenvalues.

Hence

\[
C_L
\]
increases.
\end{proof}

\subsection{Probe Stability}

\begin{definition}[Probe Stability]
A probe is stable at x if

\[
\exists \epsilon>0
\]
such that

\[
B_\epsilon(x)
\subseteq L.
\]
\end{definition}
\begin{theorem}
Stable probes occur precisely in the interior

\[
L^\circ.
\]
\end{theorem}
\begin{proof}
Interior points admit neighborhoods contained entirely within L.

Boundary points do not.
\end{proof}

\subsection{Licensing Action Functional}

Define

\[
S_L
\int_M
\Big(
|\nabla A|^2
+
|\nabla\omega|^2
+
A|\omega|
\Big)
,dV.
\]
\end{definition}

\begin{theorem}
Critical licensing configurations satisfy

\[
\Delta A
\frac12|\omega|
\]

and

\[
\Delta\omega
\frac12A,\mathrm{sgn}(\omega).
\]
\end{theorem}

\begin{proof}
Applying the Euler-Lagrange equations to

\[
S_L
\]
yields the coupled field equations.
\end{proof}

\subsection{NPI Probe Theorem}

\begin{theorem}[NPI Probe Theorem]
NPIs function as geometric probes of admissible orientation-reversing regions.

Specifically,

\[
\Pi(x)=1
\]
iff

\[
x\in L.
\]
Consequently NPIs detect the intersection

\[
L
\mathcal A_1
\cap
\Omega^-,
\]

where admissibility and orientation reversal coexist.
\end{theorem}

\begin{proof}
Directly from the definitions of licensing and the characteristic probe function.
\end{proof}

\subsection{Geometric Licensing Principle}

\begin{theorem}[Geometric Licensing Principle]
Licensing is not fundamentally a lexical property.

Licensing is membership in a negatively oriented admissible manifold.

NPIs therefore reveal the geometry of admissibility boundaries rather than merely the presence of syntactic negation.
\end{theorem}

\begin{proof}
Every licensing condition depends only on

\[
A(x)
\]
and

\[
\omega(x).
\]
No reference to lexical negation is required.

Therefore licensing is geometric.
\end{proof}

\section{Neural Dissociation Constraints and Inferential Field Operators}

\subsection{Neural State Spaces}

Let

\[
\mathcal H
\]
be a separable Hilbert space of neural states.

A neural configuration is represented by

\[
\psi \in \mathcal H.
\]
Inner products are denoted

\[
\langle \psi,\phi\rangle.
\]
Norms are

\[
|\psi|
\sqrt{\langle\psi,\psi\rangle}.
\]

\begin{definition}[Neural Representation Map]
A neural representation is a map

\[
\rho:X\rightarrow\mathcal H.
\]
\end{definition}
\begin{definition}[Population State]
The state of a neural population at time t is

\[
\psi_t.
\]
\end{definition}
\subsection{Operator Decomposition}

Define four operators

\[
C,
\quad
O,
\quad
N,
\quad
V.
\]
\begin{definition}[Parsing Operator]
\[
C:\mathcal H\rightarrow\mathcal H
\]
extracts compositional structure.
\end{definition}

\begin{definition}[Orientation Operator]
\[
O:\mathcal H\rightarrow\mathcal H
\]
constructs inferential orientation.
\end{definition}

\begin{definition}[Negation Operator]
\[
N:\mathcal H\rightarrow\mathcal H
\]
implements orientation reversal.
\end{definition}

\begin{definition}[Verification Operator]
\[
V:\mathcal H\rightarrow\mathcal H
\]
computes consistency between representation and scenario.
\end{definition}

The complete neural computation is

\[
\Psi
VNOC.
\]

\subsection{Sequential Architecture}

Given input state

\[
\psi_0,
\]
define

\[
\psi_1
C\psi_0,
\]

\[
\psi_2
O\psi_1,
\]

\[
\psi_3
N\psi_2,
\]

\[
\psi_4
V\psi_3.
\]

Thus

\[
\psi_4
VNOC\psi_0.
\]

\subsection{Functional Independence}

\begin{definition}[Operator Independence]
Operators

\[
A,B
\]
are independent iff

\[
A\neq f(B)
\]
for every measurable function f.
\]
\end{definition}

\begin{definition}[Neural Dissociation]
Two operators are neurally dissociable if there exists a perturbation

\[
P
\]
such that

\[
PA\neq A
\]
while

\[
PB=B.
\]
\end{definition}
\begin{theorem}[Dissociation Criterion]
Suppose

\[
A
\]
and

\[
B
\]
are dissociable.

Then they cannot be identical.
\end{theorem}

\begin{proof}
If

\[
A=B,
\]
then

\[
PA=PB.
\]
Contradiction.
\end{proof}

\subsection{Orientation Construction}

Let

\[
\omega
\]
denote inferential orientation.

\begin{definition}[Orientation Map]
\[
\Omega:\mathcal H\rightarrow\mathbb R.
\]
\end{definition}
Define

\[
\omega
\Omega(\psi).
\]

\begin{definition}[Orientation Population]
A neural population

\[
P_\omega
\]
encodes orientation iff

\[
\Omega(P_\omega)\neq0.
\]
\end{definition}
\subsection{Negation Dynamics}

Assume

\[
N^2=I.
\]
\begin{theorem}[Involution Property]
Every eigenvalue of N satisfies

\[
\lambda=\pm1.
\]
\end{theorem}
\begin{proof}
Let

\[
N\phi=\lambda\phi.
\]
Applying N,

\[
N^2\phi
\lambda^2\phi.
\]

Since

\[
N^2=I,
\]
\[
\lambda^2=1.
\]
Thus

\[
\lambda=\pm1.
\]
\end{proof}
\begin{definition}[Orientation-Reversal Subspace]
\[
\mathcal H_-
{
\phi:
N\phi=-\phi
}.
\]
\end{definition}

\begin{definition}[Orientation-Preserving Subspace]
\[
\mathcal H_+
{
\phi:
N\phi=\phi
}.
\]
\end{definition}

\begin{theorem}
\[
\mathcal H
\mathcal H_+
\oplus
\mathcal H_-.
\]
\end{theorem}

\begin{proof}
Spectral decomposition of involutions.
\end{proof}

\subsection{Verification Geometry}

Let

\[
S
\]
denote a scenario representation.

\begin{definition}[Verification Functional]
\[
\chi(\psi,S)
\langle\psi,S\rangle.
\]
\end{definition}

\begin{definition}[Verification Error]
\[
E_V
|\psi-S|^2.
\]
\end{definition}

\begin{theorem}
Verification succeeds iff

\[
E_V=0.
\]
\end{theorem}
\begin{proof}
Norms vanish only at equality.
\end{proof}

\subsection{Operator Non-Commutativity}

\begin{definition}
The commutator of operators A,B is

\[
[A,B]
AB-BA.
\]
\end{definition}

\begin{theorem}
If

\[
[C,N]\neq0,
\]
then parsing and negation cannot be collapsed into a single stage.
\end{theorem}

\begin{proof}
Assume a single-stage operator

\[
T.
\]
Then

\[
CN=NC.
\]
Contradiction.
\end{proof}

\begin{theorem}
If

\[
[O,N]\neq0,
\]
orientation construction and negation are computationally distinct.
\end{theorem}

\begin{proof}
Identical.
\end{proof}

\begin{theorem}
If

\[
[V,N]\neq0,
\]
verification cannot precede orientation reversal.
\end{theorem}

\begin{proof}
Order affects output.

Therefore stages are distinct.
\end{proof}

\subsection{Perturbation Theory}

Let

\[
\epsilon P
\]
be a neural perturbation.

Define

\[
N_\epsilon
N+\epsilon P.
\]

\begin{theorem}
Eigenvalues satisfy

\[
\lambda_i(\epsilon)
\lambda_i
+
\epsilon
\langle\phi_i,P\phi_i\rangle
+
O(\epsilon^2).
\]
\end{theorem}

\begin{proof}
Standard first-order perturbation theory.
\end{proof}

\subsection{Identifiability}

\begin{definition}[Operator Identifiability]
An operator A is identifiable iff

\[
A\psi=B\psi
\]
for all \psi

implies

\[
A=B.
\]
\end{definition}
\begin{theorem}
Distinct neural operators are identifiable whenever the representation space spans \mathcal H.
\end{theorem}

\begin{proof}
Operator equality on a spanning set implies equality everywhere.
\end{proof}

\subsection{Minimal Neural Architecture}

Define

\[
\mathfrak A
{C,O,N,V}.
\]

\begin{definition}[Minimal Architecture]
\[
\mathfrak A
\]
is minimal iff no proper subset computes

\[
VNOC.
\]
\end{definition}
\begin{theorem}[Minimality Theorem]
If

\[
[C,O]\neq0,
\]
\[
[C,N]\neq0,
\]
\[
[O,N]\neq0,
\]
and

\[
[V,N]\neq0,
\]
then

\[
\mathfrak A
\]
is minimal.
\end{theorem}

\begin{proof}
Removing any operator changes the resulting composition.

Therefore no proper subset reproduces

\[
VNOC.
\]
\end{proof}
\subsection{Neural Dissociation Theorem}

\begin{theorem}[Neural Dissociation Theorem]
Suppose there exist neural populations

\[
P_C,
\quad
P_O,
\quad
P_N,
\quad
P_V
\]
such that selective perturbations affect exactly one operator among

\[
C,O,N,V.
\]
Then parsing, orientation construction, negation, and verification are computationally distinct processes.
\end{theorem}

\begin{proof}
Selective perturbation establishes pairwise dissociation.

By the Dissociation Criterion, dissociated operators cannot be identical.

Therefore

\[
C\neq O,
\qquad
C\neq N,
\qquad
C\neq V,
\]
\[
O\neq N,
\qquad
O\neq V,
\qquad
N\neq V.
\]
Hence the four stages are distinct.
\end{proof}

\subsection{Strong Dissociation Corollary}

\begin{corollary}
No theory identifying negation solely with syntax, solely with semantic verification, or solely with orientation construction can reproduce the full operator structure

\[
VNOC.
\]
\end{corollary}
\begin{proof}
Each reduction removes at least one dissociable operator.

Minimality then fails.
\end{proof}

\subsection{Operator Separation Principle}

\begin{theorem}[Operator Separation Principle]
If parsing, orientation construction, negation, and verification admit independent neural realizations, then inferential computation necessarily possesses at least four irreducible stages.

Consequently negation cannot be reduced to syntax, semantics, or verification alone.
\end{theorem}

\begin{proof}
Follows from operator identifiability, non-commutativity, and neural dissociation.
\end{proof}

\section{Inferential Field Semantics and Classical Logic}

\subsection{Inferential Fields}

Let

\[
\mathcal F
(X,R,\omega,A)
\]

be an inferential field where

\[
X
\]
is a distinction manifold,

\[
R
\subseteq
X\times X
\]
is a reachability relation,

\[
\omega:X\rightarrow\mathbb R
\]
is an orientation field,

and

\[
A:X\rightarrow[0,1]
\]
is an admissibility field.

\begin{definition}[Verification Functional]
A proposition P is verified at x\in X iff

\[
\chi(P,x)=1.
\]
\end{definition}
\begin{definition}[Truth]
\[
\operatorname{True}(P,x)
\]
iff

\[
A(x)=1
\]
and

\[
\chi(P,x)=1.
\]
\end{definition}
Thus truth is admissible verification.

\subsection{Logical Interpretation}

Let

\[
\mathcal L
\]
denote classical propositional logic.

Define

\[
\mathcal T
:
\mathcal L
\rightarrow
\mathcal F
\]
by

\[
P
\mapsto
{x:\chi(P,x)=1}.
\]
\begin{definition}[Interpretation Functor]
The map

\[
\mathcal T
\]
is called the inferential interpretation functor.
\end{definition}

\subsection{Conjunction}

\begin{definition}
For propositions

\[
P,Q,
\]
define

\[
P\wedge Q
\]
by

\[
\chi(P\wedge Q,x)
\chi(P,x)\chi(Q,x).
\]
\end{definition}

\begin{theorem}
\[
\mathcal T(P\wedge Q)
\mathcal T(P)
\cap
\mathcal T(Q).
\]
\end{theorem}

\begin{proof}
\[
\chi(P\wedge Q,x)=1
\]
iff

\[
\chi(P,x)=1
\]
and

\[
\chi(Q,x)=1.
\]
Hence

\[
x
\in
\mathcal T(P)\cap\mathcal T(Q).
\]
\end{proof}
\subsection{Disjunction}

\begin{definition}
\[
\chi(P\vee Q,x)
\max
{
\chi(P,x),
\chi(Q,x)
}.
\]
\end{definition}

\begin{theorem}
\[
\mathcal T(P\vee Q)
\mathcal T(P)
\cup
\mathcal T(Q).
\]
\end{theorem}

\begin{proof}
Immediate.
\end{proof}

\subsection{Implication}

\begin{definition}[Reachability Implication]
\[
P\Rightarrow Q
\]
iff

\[
\mathcal T(P)
\subseteq
\mathcal T(Q).
\]
\end{definition}
\begin{definition}[Inferential Entailment]
\[
P\models_R Q
\]
iff

\[
R(x,y)
\]
and

\[
\chi(P,x)=1
\]
imply

\[
\chi(Q,y)=1.
\]
\end{definition}
\begin{theorem}
Inferential entailment is transitive.
\end{theorem}

\begin{proof}
Suppose

\[
P\models_R Q
\]
and

\[
Q\models_R S.
\]
Then

\[
P
\rightarrow
Q
\rightarrow
S.
\]
Hence

\[
P\models_R S.
\]
\end{proof}
\subsection{Negation}

Let

\[
N:X\rightarrow X
\]
be orientation reversal.

\[
N^2=I.
\]
\begin{definition}[Inferential Negation]
\[
\chi(\neg P,x)
\chi(P,N(x)).
\]
\end{definition}

\begin{theorem}
Negation corresponds to orientation reversal.
\end{theorem}

\begin{proof}
By definition,

\[
\neg P
\]
is evaluated at

\[
N(x).
\]
Therefore truth evaluation occurs after orientation reversal.
\end{proof}

\subsection{Double Negation}

\begin{theorem}[Double Negation Restoration]
\[
\chi(\neg\neg P,x)
\chi(P,x).
\]
\end{theorem}

\begin{proof}
\[
\chi(\neg\neg P,x)
\chi(\neg P,N(x))

\chi(P,N^2(x)).
\]

Since

\[
N^2=I,
\]
\[
\chi(\neg\neg P,x)
\chi(P,x).
\]
\end{proof}

\begin{corollary}
\[
\neg\neg P
\equiv
P.
\]
\end{corollary}
\subsection{Truth Regions}

Define

\[
T(P)
{
x:
\operatorname{True}(P,x)
}.
\]

\begin{theorem}
\[
T(P)
\mathcal A_1
\cap
\mathcal T(P).
\]
\end{theorem}

\begin{proof}
Truth requires both admissibility and verification.
\end{proof}

\subsection{Soundness}

\begin{theorem}[Soundness]
If

\[
\vdash P
\]
in classical propositional logic, then

\[
\models_{\mathcal F} P.
\]
\end{theorem}
\begin{proof}
Each logical connective preserves its set-theoretic interpretation.

Classical derivations therefore preserve truth regions.

Hence every theorem is valid in inferential fields.
\end{proof}

\subsection{Completeness}

\begin{theorem}[Completeness]
If

\[
\models_{\mathcal F} P,
\]
then

\[
\vdash P.
\]
\end{theorem}
\begin{proof}
Truth regions form a Boolean algebra under

\[
\cap,
\cup,
\text{and }
N.
\]
Stone representation implies equivalence with classical valuations.

Therefore every inferentially valid proposition is classically derivable.
\end{proof}

\subsection{Boolean Structure}

\begin{theorem}
The collection

\[
\mathfrak B
{
T(P):P\in\mathcal L
}
\]

forms a Boolean algebra.
\end{theorem}

\begin{proof}
Closure under

\[
\cap,
\cup,
N
\]
follows from conjunction, disjunction, and negation.

Classical Boolean identities are inherited.
\end{proof}

\subsection{Orientation-Neutral Regions}

Define

\[
B
{
x:
\omega(x)=0
}.
\]

\begin{definition}[Boundary Region]
B is the orientation-neutral boundary.
\end{definition}

\begin{theorem}
Classical bivalence may fail on B.
\end{theorem}

\begin{proof}
Orientation reversal becomes undefined or degenerate at

\[
\omega=0.
\]
Hence truth-value assignment need not be unique.
\end{proof}

\subsection{Intuitionistic Boundary Logic}

\begin{definition}
A proposition is stable iff

\[
\omega(x)\neq0
\]
throughout its truth region.
\]
\end{definition}

\begin{theorem}
Boundary regions support intuitionistic semantics.
\end{theorem}

\begin{proof}
Double-negation elimination requires nondegenerate orientation.

At

\[
\omega=0,
\]
only

\[
P\rightarrow\neg\neg P
\]
is guaranteed.
\end{proof}

\subsection{Modal Reachability}

Introduce operators

\[
\Box,
\quad
\Diamond.
\]
\begin{definition}
\[
\Box P
\]
iff

\[
\forall y,
,
R(x,y)
\Rightarrow
\chi(P,y)=1.
\]
\end{definition}
\begin{definition}
\[
\Diamond P
\]
iff

\[
\exists y,
,
R(x,y)
\wedge
\chi(P,y)=1.
\]
\end{definition}
\begin{theorem}
\[
\Box P
\Rightarrow
\Diamond P.
\]
\end{theorem}
\begin{proof}
Reflexivity of R.
\end{proof}

\subsection{Reachability Frames}

\begin{definition}
The frame associated with an inferential field is

\[
(X,R).
\]
\end{definition}
\begin{theorem}
Every inferential field determines a Kripke frame.
\end{theorem}

\begin{proof}
The reachability relation acts as accessibility.
\end{proof}

\begin{theorem}
Every Kripke frame can be embedded into an inferential field.
\end{theorem}

\begin{proof}
Assign

\[
\omega\equiv1,
\qquad
A\equiv1.
\]
Then the remaining structure is exactly

\[
(X,R).
\]
\end{proof}
\subsection{Representation Theorem}

\begin{theorem}[Representation Theorem]
Classical propositional logic embeds faithfully into inferential-field semantics via

\[
\mathcal T.
\]
Logical connectives become geometric operations on admissible verification regions, while negation becomes orientation reversal.
\end{theorem}

\begin{proof}
The preceding conjunction, disjunction, implication, and negation theorems establish preservation of logical structure.

Faithfulness follows from soundness and completeness.
\end{proof}

\subsection{Negation Before Logic Theorem}

\begin{theorem}[Negation Before Logic]
Let

\[
\mathcal F
(X,R,\omega,A).
\]

Then logical negation is derivable from the existence of an involutive orientation-reversal map

\[
N:X\rightarrow X.
\]
Consequently orientation precedes logical negation, and admissible reachability precedes truth.
\end{theorem}

\begin{proof}
Truth evaluation requires verification.

Verification requires admissibility.

Negation acts on orientation prior to verification.

Therefore

\[
A
;\prec;
\omega
;\prec;
\neg
;\prec;
\text{truth}.
\]
Hence negation is geometrically prior to logic.
\end{proof}

\section{Distinguishability Geometry and Inferential Displacement}

\subsection{Distinguishability Spaces}

Let

\[
X
\]
be a nonempty set.

\begin{definition}[Distinguishability Relation]
A distinguishability relation is a symmetric relation

\[
\not\sim
;\subseteq;
X\times X.
\]
\end{definition}
Its complement

\[
\sim
\]
is the indistinguishability relation.

\begin{definition}[Distinguishability Space]
A distinguishability space is the pair

\[
\mathcal D
(X,\sim).
\]
\end{definition}

\begin{definition}[Equivalence Class]
The equivalence class of x is

\[
[x]
{y\in X:y\sim x}.
\]
\end{definition}

\begin{definition}[Quotient Space]
\[
X/{\sim}
{
[x]:
x\in X
}.
\]
\end{definition}

\subsection{Distinguishability Metrics}

Let

\[
d_D
\]
be a distinguishability metric.

\begin{definition}[Distinguishability Metric]
\[
d_D:X\times X\rightarrow\mathbb R_{\ge0}
\]
satisfies

\[
d_D(x,y)=0
\iff
x\sim y.
\]
\end{definition}
\begin{definition}[Distinguishability Neighborhood]
\[
B_r(x)
{
y:
d_D(x,y)<r
}.
\]
\end{definition}

\begin{definition}[Distinguishability Volume]
\[
V_D(x,r)
\mu(B_r(x)).
\]
\end{definition}

\subsection{Reachability Embedding}

Let

\[
(X,R)
\]
be a reachability space.

\begin{definition}[Reachability-Induced Distinguishability]
\[
d_D(x,y)
|d_R(x,y)-d_R(y,x)|.
\]
\end{definition}

\begin{theorem}
Every metric reachability space induces a distinguishability geometry.
\end{theorem}

\begin{proof}
The induced metric separates points according to asymmetry of reachability.

Thus distinguishability emerges from reachability structure.
\end{proof}

\subsection{Distinguishability Deficit}

Let

\[
L_{\mathrm{opt}}(x)
\]
denote optimal reachable distinguishability.

Let

\[
L_{\mathrm{act}}(x)
\]
denote actual distinguishability.

\begin{definition}[Distinguishability Deficit]
\[
\delta_D(x)
L_{\mathrm{opt}}(x)

L_{\mathrm{act}}(x).
\]
\end{definition}

\begin{definition}[Global Deficit]
\[
\Delta_D
\int_X
\delta_D(x)
,d\mu(x).
\]
\end{definition}

\begin{theorem}
\[
\Delta_D\ge0.
\]
\end{theorem}
\begin{proof}
Optimal distinguishability dominates actual distinguishability.

Hence

\[
\delta_D(x)\ge0.
\]
Integrating preserves positivity.
\end{proof}

\subsection{Negation as Distinguishability Displacement}

Let

\[
N:X\rightarrow X
\]
be orientation reversal.

\begin{definition}[Negation Displacement]
\[
\eta_N(x)
d_D(x,N(x)).
\]
\]
\end{definition}

\begin{definition}[Average Negation Displacement]
\[
\bar\eta_N
\int_X
\eta_N(x)
,d\mu(x).
\]
\end{definition}

\begin{theorem}
Negation distortion is a special case of distinguishability displacement.
\end{theorem}

\begin{proof}
By definition,

\[
D_N(x)
d_R(x,N(x)).
\]

Since

\[
d_D
\]
is induced by

\[
d_R,
\]
negation distortion becomes a distinguishability displacement.
\end{proof}

\subsection{Distinguishability Tensor Geometry}

Assume X is a smooth manifold.

Let

\[
g_{ij}
\]
denote the distinguishability metric tensor.

\begin{definition}[Deficit Tensor]
\[
\Theta_{ij}
g^{\mathrm{opt}}_{ij}

g^{\mathrm{act}}_{ij}.
\]
\end{definition}

\begin{definition}[Tensor Norm]
\[
|\Theta|^2
\Theta_{ij}\Theta^{ij}.
\]
\end{definition}

\begin{theorem}
\[
\Theta=0
\]
iff distinguishability is optimal everywhere.
\end{theorem}

\begin{proof}
Vanishing tensor implies

\[
g^{\mathrm{opt}}
g^{\mathrm{act}}.
\]
\end{proof}

\subsection{Distinguishability Curvature}

Let

\[
\nabla
\]
be the Levi-Civita connection.

\begin{definition}[Distinguishability Curvature]
\[
R^i{}_{jkl}
\]
is the Riemann curvature tensor associated with

\[
g_{ij}.
\]
\end{definition}
\begin{definition}[Scalar Distinguishability Curvature]
\[
K_D
g^{ij}R_{ij}.
\]
\end{definition}

\begin{theorem}
Regions of large

\[
|K_D|
\]
exhibit rapid distinguishability deformation.
\end{theorem}

\begin{proof}
Curvature measures local deviation from Euclidean distinguishability structure.
\end{proof}

\subsection{Admissibility Quotients}

Let

\[
A:X\rightarrow[0,1]
\]
be admissibility.

\begin{definition}[Admissibility Equivalence]
\[
x\sim_A y
\]
iff

\[
A(x)=A(y).
\]
\end{definition}
\begin{definition}[Admissibility Quotient]
\[
X/{\sim_A}.
\]
\end{definition}
\begin{theorem}
The admissibility quotient partitions distinguishability space into admissibility strata.
\end{theorem}

\begin{proof}
Equality of admissibility is an equivalence relation.
\end{proof}

\subsection{Distinguishability Entropy}

\begin{definition}[Local Distinguishability Entropy]
\[
S_D(x)
\log V_D(x,r).
\]
\end{definition}

\begin{definition}[Global Distinguishability Entropy]
\[
S_D
\int_X
S_D(x)
,d\mu(x).
\]
\end{definition}

\begin{theorem}
Increasing distinguishability volume increases entropy.
\end{theorem}

\begin{proof}
The logarithm is monotone.
\end{proof}

\subsection{Repair Operators}

Let

\[
\mathcal R
:
X\rightarrow X
\]
be a repair operator.

\begin{definition}[Repair Condition]
\[
\delta_D(\mathcal R(x))
<
\delta_D(x).
\]
\end{definition}
\begin{definition}[Perfect Repair]
\[
\delta_D(\mathcal R(x))
0. 

\]
\end{definition}

\begin{theorem}
Repeated repair generates a monotone deficit sequence

\[
\delta_D^{(n+1)}
\le
\delta_D^{(n)}.
\]
\end{theorem}
\begin{proof}
Repair reduces deficit by definition.
\end{proof}

\subsection{Continuation Geometry}

Let

\[
C_t
\]
denote continuation dynamics.

\begin{definition}[Continuation Functional]
\[
\Gamma(x)
\int_0^\infty
e^{-\lambda t}
,
L(C_t(x))
,dt.
\]
\end{definition}

\begin{definition}[Continuation Capacity]
\[
\Phi_C
\int_X
\Gamma(x)
,d\mu(x).
\]
\end{definition}

\begin{theorem}
Repair increases continuation capacity.
\end{theorem}

\begin{proof}
Repair restores reachable distinguishability.

Therefore future trajectories increase.

Hence

\[
\Phi_C
\]
increases.
\end{proof}

\subsection{Admissibility Emergence}

\begin{definition}[Admissibility Functional]
\[
\Lambda(x)
f
\big(
\delta_D(x),
\Gamma(x)
\big).
\]
\end{definition}

\begin{theorem}
Admissibility increases with continuation and decreases with deficit.
\end{theorem}

\begin{proof}
By monotonicity of f.
\end{proof}

\subsection{Hierarchy Theorem}

\begin{theorem}[Distinction Hierarchy]
For every admissible system,

\[
\text{Distinction}
\rightarrow
\text{Information}
\rightarrow
\text{Entropy}
\rightarrow
\text{Repair}
\rightarrow
\text{Continuation}
\rightarrow
\text{Admissibility}.
\]
\end{theorem}
\begin{proof}
Distinctions generate informational differences.

Information determines distinguishability entropy.

Entropy generates deficit.

Repair reduces deficit.

Reduced deficit increases continuation.

Continuation determines admissibility.
\end{proof}

\subsection{Unification Theorem}

\begin{theorem}[Unification Theorem]
Negation distortion, admissibility geometry, repair theory, continuation theory, and distinguishability geometry are all projections of a common underlying displacement structure.

Specifically,

\[
D_N
\subseteq
\eta_N
\subseteq
\delta_D
\subseteq
\Theta
\subseteq
\Lambda.
\]
Thus inferential negation is a special case of distinguishability displacement, which is itself a special case of admissibility deformation.
\end{theorem}

\begin{proof}
Each quantity is obtained by successive abstraction from reachability displacement to distinguishability deficit to admissibility structure.

Therefore the inclusions follow by construction.
\end{proof}

\subsection{Fundamental Geometric Principle}

\begin{theorem}[Fundamental Geometric Principle]
Objects, propositions, memories, repairs, and admissible futures are not primitive entities.

They are stable regions of distinguishability geometry whose persistence depends on the maintenance of recoverable distinctions.

Consequently admissibility is the terminal invariant of distinction-preserving dynamics.
\end{theorem}

\begin{proof}
Distinguishability generates information.

Information generates entropy.

Entropy necessitates repair.

Repair enables continuation.

Continuation defines admissibility.

Hence admissibility emerges as the terminal invariant of the hierarchy.
\end{proof}

\section{Geometric Foundations of Admissibility}

\subsection{Fundamental Configuration Space}

Let

\[
M
\]
be a smooth differentiable manifold.

Define the fundamental field tuple

\[
\Psi
(D,\omega,\mathcal M,\mathcal R,\Lambda)
\]

where

\[
D:M\rightarrow\mathbb R
\]
is the distinguishability field,

\[
\omega:M\rightarrow\mathbb R
\]
is the orientation field,

\[
\mathcal M:M\rightarrow\mathbb R
\]
is the memory field,

\[
\mathcal R:M\rightarrow\mathbb R
\]
is the repair field,

and

\[
\Lambda:M\rightarrow\mathbb R
\]
is the admissibility field.

\subsection{Fundamental Principle}

\begin{axiom}[Recoverable Distinction Principle]
A configuration is physically, cognitively, and inferentially meaningful only to the extent that distinctions remain recoverable under admissible transformations.
\end{axiom}

\begin{axiom}[Continuation Principle]
Admissibility is determined not by present state alone but by the existence of recoverable future trajectories.
\end{axiom}

\subsection{Fundamental Action}

Define

\[
S[\Psi]
\int_M
\mathcal L
,dV.
\]

The Lagrangian density is

\[
\mathcal L
\mathcal L_D
+
\mathcal L_\omega
+
\mathcal L_M
+
\mathcal L_R
+
\mathcal L_\Lambda
+
\mathcal L_{\mathrm{int}}.
\]

\subsection{Distinguishability Sector}

\[
\mathcal L_D
\frac12
|\nabla D|^2

V_D(D).
\]

Define

\[
V_D(D)
\frac{\alpha}{2}D^2
+
\frac{\beta}{4}D^4.
\]

\begin{theorem}
Critical points satisfy

\[
\Delta D
\alpha D
+
\beta D^3.
\]
\end{theorem}

\begin{proof}
Euler-Lagrange variation.
\end{proof}

\subsection{Orientation Sector}

\[
\mathcal L_\omega
\frac12
|\nabla\omega|^2

V_\omega(\omega).
\]

Choose

\[
V_\omega
\frac{\gamma}{2}\omega^2.
\]

Then

\[
\Delta\omega
\gamma\omega.
\]

\subsection{Memory Sector}

\[
\mathcal L_M
\frac12
|\nabla\mathcal M|^2

\frac{\mu}{2}\mathcal M^2.
\]

Hence

\[
\Delta\mathcal M
\mu\mathcal M.
\]

\begin{definition}[Persistence Density]
\[
P
\mathcal M^2.
\]
\end{definition}

\subsection{Repair Sector}

\[
\mathcal L_R
\frac12
|\nabla\mathcal R|^2

V_R(\mathcal R).
\]

Let

\[
V_R
\frac{\rho}{2}\mathcal R^2.
\]

Then

\[
\Delta\mathcal R
\rho\mathcal R.
\]

\begin{definition}[Repair Capacity]
\[
\Phi_R
\int_M
\mathcal R
,dV.
\]
\end{definition}

\subsection{Admissibility Sector}

\[
\mathcal L_\Lambda
\frac12
|\nabla\Lambda|^2

V_\Lambda(\Lambda).
\]

Choose

\[
V_\Lambda
\frac{\lambda}{2}\Lambda^2.
\]

Then

\[
\Delta\Lambda
\lambda\Lambda.
\]

\subsection{Interaction Terms}

Define

\[
\mathcal L_{\mathrm{int}}
aD\omega
+
bD\mathcal M
+
cD\mathcal R
+
d\mathcal M\mathcal R
+
e\mathcal R\Lambda.
\]

\subsection{Coupled Field Equations}

Variation yields

\[
\Delta D
\alpha D
+
\beta D^3
+
a\omega
+
b\mathcal M
+
c\mathcal R,
\]

\[
\Delta\omega
\gamma\omega
+
aD,
\]

\[
\Delta\mathcal M
\mu\mathcal M
+
bD
+
d\mathcal R,
\]

\[
\Delta\mathcal R
\rho\mathcal R
+
cD
+
d\mathcal M
+
e\Lambda,
\]

\[
\Delta\Lambda
\lambda\Lambda
+
e\mathcal R.
\]

\subsection{Fundamental Dependency Chain}

The coupling structure induces

\[
D
\rightarrow
\omega
\rightarrow
\mathcal M
\rightarrow
\mathcal R
\rightarrow
\Lambda.
\]
\begin{theorem}
Admissibility depends indirectly on distinguishability through the entire chain.
\end{theorem}

\begin{proof}
The coupled equations show

\[
\Lambda
\]
depends on

\[
\mathcal R,
\]
which depends on

\[
\mathcal M,
\]
which depends on

\[
D.
\]
\end{proof}
\subsection{Distinguishability Current}

Define

\[
J_D
\nabla D.
\]

\begin{definition}[Distinction Flux]
\[
\Phi_D
\int_{\partial M}
J_D\cdot n
,dS.
\]
\end{definition}

\begin{theorem}
\[
\Phi_D
\int_M
\Delta D
,dV.
\]
\end{theorem}

\begin{proof}
Divergence theorem.
\end{proof}

\subsection{Orientation Current}

\[
J_\omega
\nabla\omega.
\]

\begin{definition}[Orientation Flux]
\[
\Phi_\omega
\int_{\partial M}
J_\omega\cdot n
,dS.
\]
\end{definition}

\subsection{Memory Current}

\[
J_M
\nabla\mathcal M.
\]

\begin{definition}[Persistence Flux]
\[
\Phi_M
\int_{\partial M}
J_M\cdot n
,dS.
\]
\end{definition}

\subsection{Repair Current}

\[
J_R
\nabla\mathcal R.
\]

\begin{definition}[Repair Flux]
\[
\Phi_R
\int_{\partial M}
J_R\cdot n
,dS.
\]
\end{definition}

\subsection{Admissibility Current}

\[
J_\Lambda
\nabla\Lambda.
\]

\begin{definition}[Admissibility Flux]
\[
\Phi_\Lambda
\int_{\partial M}
J_\Lambda\cdot n
,dS.
\]
\end{definition}

\subsection{Conservation Law}

Define total geometric charge

\[
Q
\int_M
(D+\omega+\mathcal M+\mathcal R+\Lambda)
,dV.
\]

\begin{theorem}
If interaction terms are symmetric,

\[
\frac{dQ}{dt}=0.
\]
\end{theorem}
\begin{proof}
Integrating the coupled field equations causes interaction contributions to cancel pairwise.
\end{proof}

\subsection{Recoverability Functional}

Define

\[
\mathfrak R
\int_M
D\mathcal M
,dV.
\]

\begin{definition}[Recoverability]
A system is recoverable iff

\[
\mathfrak R>0.
\]
\end{definition}
\subsection{Continuation Functional}

Define

\[
\mathfrak C
\int_M
\mathcal M\mathcal R
,dV.
\]

\begin{definition}[Continuation Capacity]
\[
\Phi_C
\mathfrak C.
\]
\end{definition}

\begin{theorem}
Repair without memory yields

\[
\Phi_C=0.
\]
\end{theorem}
\begin{proof}
If

\[
\mathcal M=0,
\]
then

\[
\mathfrak C=0.
\]
\end{proof}
\subsection{Admissibility Functional}

Define

\[
\mathfrak A
\int_M
\Lambda
\mathcal R
\mathcal M
,dV.
\]

\begin{definition}[Global Admissibility]
\[
A_{\mathrm{global}}
\mathfrak A.
\]
\end{definition}

\subsection{Entropy as Deficit}

Define

\[
E
D_{\mathrm{opt}}

D.
\]

\begin{definition}[Distinguishability Entropy]
\[
S_D
\int_M
E
,dV.
\]
\end{definition}

\begin{theorem}
Repair reduces entropy deficit.
\end{theorem}

\begin{proof}
Repair increases realized distinguishability.

Hence

\[
E
\]
decreases.
\end{proof}

\subsection{Master Stability Functional}

Define

\[
\Sigma
\int_M
\left(
D^2
+
\omega^2
+
\mathcal M^2
+
\mathcal R^2
+
\Lambda^2
\right)
dV.
\]

\begin{theorem}
Stable systems are local minima of

\[
\Sigma.
\]
\end{theorem}
\begin{proof}
Standard Lyapunov argument.
\end{proof}

\subsection{Grand Unification Theorem}

\begin{theorem}[Grand Unification]
Distinction, orientation, memory, repair, continuation, and admissibility are not independent primitives.

They arise as coupled fields governed by a single variational principle

\[
\delta S=0.
\]
All previously defined structures—

\[
D_N,
\quad
\delta_D,
\quad
\Theta,
\quad
\Gamma,
\quad
\Lambda,
\]
appear as derived quantities of the unified field configuration

\[
\Psi
(D,\omega,\mathcal M,\mathcal R,\Lambda).
\]
\end{theorem}

\begin{proof}
Each quantity is either a field, a functional of fields, or a derived invariant of the action.

Therefore all emerge from the stationary points of

\[
S.
\]
\end{proof}
\subsection{The Fate of Distinguishability}

\begin{theorem}[Terminal Principle]
A system persists precisely to the extent that distinguishability remains recoverable.

Memory stores recoverable distinctions.

Repair restores lost distinctions.

Continuation propagates distinctions.

Admissibility selects distinctions capable of persistence.

Consequently,

\[
\text{Admissibility}
\text{Recoverable Future Distinguishability}.
\]
\end{theorem}

\begin{proof}
Immediate from the dependency chain

\[
D
\rightarrow
\omega
\rightarrow
\mathcal M
\rightarrow
\mathcal R
\rightarrow
\Lambda.
\]

The terminal invariant of the chain is admissibility, and every preceding stage exists only insofar as distinctions remain recoverable.
\end{proof}

\newpage

\begin{thebibliography}{99}

\bibitem{Agmon2019}
Agmon, G., Loewenstein, Y., & Grodzinsky, Y. (2019).
Measuring the cognitive cost of downward monotonicity by controlling for negative polarity.
\textit{Glossa: A Journal of General Linguistics}, 4(1), 36.

\bibitem{Agmon2022}
Agmon, G., Loewenstein, Y., & Grodzinsky, Y. (2022).
Negative sentences exhibit a sustained effect in delayed verification tasks.
\textit{Journal of Experimental Psychology: Learning, Memory, and Cognition}, 48(1), 122--141.

\bibitem{Amunts2020}
Amunts, K., Mohlberg, H., Bludau, S., & Zilles, K. (2020).
Julich-Brain: A 3D probabilistic atlas of the human brain's cytoarchitecture.
\textit{Science}, 369(6506), 988--992.

\bibitem{Baker1970}
Baker, C. L. (1970).
Double negatives.
\textit{Linguistic Inquiry}, 1, 169--186.

\bibitem{BarwiseCooper1981}
Barwise, J., & Cooper, R. (1981).
Generalized quantifiers and natural language.
\textit{Linguistics and Philosophy}, 4, 159--219.

\bibitem{Broca1861}
Broca, P. (1861).
Remarques sur le siège de la faculté du langage articulé.
\textit{Bulletins de la Société Anatomique de Paris}, 6, 330--357.

\bibitem{CaramazzaZurif1976}
Caramazza, A., & Zurif, E. (1976).
Dissociation of algorithmic and heuristic processes in language comprehension.
\textit{Brain and Language}, 3, 572--582.

\bibitem{CarpenterJust1975}
Carpenter, P. A., & Just, M. A. (1975).
Sentence comprehension: A psycholinguistic processing model of verification.
\textit{Psychological Review}, 82, 45--73.

\bibitem{Chomsky1995}
Chomsky, N. (1995).
\textit{The Minimalist Program}.
Cambridge, MA: MIT Press.

\bibitem{ClarkChase1972}
Clark, H. H., & Chase, W. G. (1972).
On the process of comparing sentences against pictures.
\textit{Cognitive Psychology}, 3, 472--517.

\bibitem{Crnic2014}
Crnič, L. (2014).
Non-monotonicity in NPI licensing.
\textit{Natural Language Semantics}, 22, 169--217.

\bibitem{Dehaene2003}
Dehaene, S., Piazza, M., Pinel, P., & Cohen, L. (2003).
Three parietal circuits for number processing.
\textit{Cognitive Neuropsychology}, 20, 487--506.

\bibitem{Deschamps2015}
Deschamps, I., Agmon, G., Loewenstein, Y., & Grodzinsky, Y. (2015).
The processing cost of downward entailingness: The representation and verification of comparative constructions.
In \textit{Proceedings of Sinn und Bedeutung 19}.

\bibitem{Fauconnier1975}
Fauconnier, G. (1975).
Polarity and the scale principle.
\textit{Chicago Linguistic Society}, 11, 188--199.

\bibitem{Fodor1983}
Fodor, J. A. (1983).
\textit{The Modularity of Mind}.
Cambridge, MA: MIT Press.

\bibitem{Frege1879}
Frege, G. (1879).
\textit{Begriffsschrift}.
Halle: Nebert.

\bibitem{Gajewski2011}
Gajewski, J. (2011).
Licensing strong NPIs.
\textit{Natural Language Semantics}, 19, 109--148.

\bibitem{Geurts2010}
Geurts, B., Katsos, N., Cummins, C., Moons, J., & Noordman, L. (2010).
Scalar quantifiers: Logic, acquisition, and processing.
\textit{Language and Cognitive Processes}, 25, 130--148.

\bibitem{Goodglass1968}
Goodglass, H. (1968).
Studies on the grammar of aphasics.
\textit{Cortex}, 4, 193--216.

\bibitem{Grodzinsky1986}
Grodzinsky, Y. (1986).
Language deficits and the theory of syntax.
\textit{Brain and Language}, 27, 135--159.

\bibitem{GrodzinskyEtAl2020}
Grodzinsky, Y., Deschamps, I., Pieperhoff, P., Iannilli, F.,
Agmon, G., Loewenstein, Y., & Amunts, K. (2020).
Logical negation mapped onto the brain.
\textit{Brain Structure and Function}, 225, 19--31.

\bibitem{GrodzinskyEtAl2021}
Grodzinsky, Y., Pieperhoff, P., & Thompson, C. K. (2021).
Stable brain loci for the processing of complex syntax.
\textit{Cortex}, 142, 252--271.

\bibitem{GrodzinskyEtAl2023}
Grodzinsky, Y., Behrent, K., Agmon, G., Bittner, N.,
Jockwitz, C., Caspers, S., Amunts, K., & Heim, S. (2023).
A linguistic complexity pattern that defies aging:
The processing of multiple negations.

\bibitem{Horn1989}
Horn, L. R. (1989).
\textit{A Natural History of Negation}.
Chicago: University of Chicago Press.

\bibitem{Homer2021}
Homer, V. (2021).
Negative polarity items.
In D. Gutzmann et al. (Eds.),
\textit{The Wiley Blackwell Companion to Semantics}.

\bibitem{Kaup2001}
Kaup, B. (2001).
Negation and its impact on the accessibility of text information.
\textit{Memory & Cognition}, 29, 960--967.

\bibitem{KaupZwaan2003}
Kaup, B., & Zwaan, R. (2003).
Effects of negation and situational presence on language comprehension.
\textit{Journal of Experimental Psychology: Learning, Memory, and Cognition}, 29, 439--446.

\bibitem{Klima1964}
Klima, E. S. (1964).
Negation in English.
In J. A. Fodor & J. Katz (Eds.),
\textit{The Structure of Language}.
Englewood Cliffs: Prentice Hall.

\bibitem{Ladusaw1980}
Ladusaw, W. (1980).
\textit{Polarity Sensitivity as Inherent Scope Relations}.
Doctoral dissertation,
University of Texas at Austin.

\bibitem{Lakoff1970}
Lakoff, G. (1970).
Linguistics and natural logic.
\textit{Synthese}, 22, 151--271.

\bibitem{Montague1974}
Montague, R. (1974).
\textit{Formal Philosophy}.
New Haven: Yale University Press.

\bibitem{Russell1903}
Russell, B. (1903).
\textit{The Principles of Mathematics}.
Cambridge: Cambridge University Press.

\bibitem{Szymanik2016}
Szymanik, J. (2016).
\textit{Quantifiers and Cognition}.
Cham: Springer.

\bibitem{Tan2023}
Tan, I.-A., Kugler-Etinger, N., & Grodzinsky, Y. (2023).
Do two negatives make a positive?
Language and logic in language processing.
\textit{Language, Cognition and Neuroscience}, 38(7), 1027--1043.

\bibitem{Tan2024}
Tan, I.-A., Segal, N., & Grodzinsky, Y. (2024).
The domains of monotonicity processing.
\textit{Journal of Semantics}, 41, 77--101.

\bibitem{Wason1961}
Wason, P. C. (1961).
Response to affirmative and negative binary statements.
\textit{British Journal of Psychology}, 52, 133--142.

\bibitem{Wason1965}
Wason, P. C. (1965).
The contexts of plausible denial.
\textit{Journal of Verbal Learning and Verbal Behavior}, 4, 7--11.

\bibitem{XiangEtAl2016}
Xiang, M., Grove, J., & Giannakidou, A. (2016).
Dependency-dependent interference:
NPI processing and intervention.
\textit{Glossa}, 1(1), 38.

\bibitem{Zurif1980}
Zurif, E. B. (1980).
Language mechanisms: A neuropsychological perspective.
\textit{American Scientist}, 68, 305--311.

\end{thebibliography}

\end{document}
